\documentclass{article}

% if you need to pass options to natbib, use, e.g.:
%     \PassOptionsToPackage{numbers, compress}{natbib}
% before loading neurips_2026

% The authors should use one of these tracks.
% Before accepting by the NeurIPS conference, select one of the options below.
% 0. "default" for submission
 % \usepackage{neurips_2026}
% the "default" option is equal to the "main" option, which is used for the Main Track with double-blind reviewing.
% 1. "main" option is used for the Main Track
%  \usepackage[main]{neurips_2026}
% 2. "position" option is used for the Position Paper Track
%  \usepackage[position]{neurips_2026}
% 3. "dandb" option is used for the Datasets & Benchmarks Track
 % \usepackage[dandb]{neurips_2026}
% 4. "creativeai" option is used for the Creative AI Track
%  \usepackage[creativeai]{neurips_2026}
% 5. "sglblindworkshop" option is used for the Workshop with single-blind reviewing
 % \usepackage[sglblindworkshop]{neurips_2026}
% 6. "dblblindworkshop" option is used for the Workshop with double-blind reviewing
%  \usepackage[dblblindworkshop]{neurips_2026}

% After being accepted, the authors should add "final" behind the track to compile a camera-ready version.
% 1. Main Track
 \usepackage[main, final]{neurips_2026}
% 2. Position Paper Track
%  \usepackage[position, final]{neurips_2026}
% 3. Datasets & Benchmarks Track
 % \usepackage[dandb, final]{neurips_2026}
% 4. Creative AI Track
%  \usepackage[creativeai, final]{neurips_2026}
% 5. Workshop with single-blind reviewing
%  \usepackage[sglblindworkshop, final]{neurips_2026}
% 6. Workshop with double-blind reviewing
%  \usepackage[dblblindworkshop, final]{neurips_2026}
% Note. For the workshop paper template, both \title{} and \workshoptitle{} are required, with the former indicating the paper title shown in the title and the latter indicating the workshop title displayed in the footnote.
% For workshops (5., 6.), the authors should add the name of the workshop, "\workshoptitle" command is used to set the workshop title.
% \workshoptitle{WORKSHOP TITLE}

% "preprint" option is used for arXiv or other preprint submissions
 % \usepackage[preprint]{neurips_2026}

% to avoid loading the natbib package, add option nonatbib:
%    \usepackage[nonatbib]{neurips_2026}

\usepackage[utf8]{inputenc} % allow utf-8 input
\usepackage[T1]{fontenc}    % use 8-bit T1 fonts
\usepackage{hyperref}       % hyperlinks
\usepackage{url}            % simple URL typesetting
\usepackage{amsmath}
\usepackage{amssymb}
\usepackage{mathtools}
\usepackage{amsthm}

\usepackage{enumitem}

\usepackage{booktabs}
\usepackage{multirow}
\usepackage{tabularx}
\usepackage{colortbl}
\usepackage{listings}
\usepackage{xcolor}
\newcommand{\cgap}[1]{\textcolor{red}{#1}}
\lstset{
  language=Python,
  basicstyle=\ttfamily\small,
  frame=single,
  breaklines=true,
  columns=fullflexible,
  keywordstyle=\color{blue},
  commentstyle=\color{gray},
  stringstyle=\color{red},
  showstringspaces=false
}



\definecolor{lightgray}{gray}{0.95}
\usepackage{tcolorbox}
\tcbuselibrary{breakable}
\usepackage{titletoc}
\usepackage{kotex}
\usepackage{subcaption}
% preamble에 추가
\usepackage{pifont}
\newcommand{\cmark}{\ding{51}}
\newcommand{\xmark}{\ding{55}}
\usepackage{verbatim}
% Note. For the workshop paper template, both \title{} and \workshoptitle{} are required, with the former indicating the paper title shown in the title and the latter indicating the workshop title displayed in the footnote. 
\title{Teaching VLMs What to Say, Not How to Reason: \\Rethinking Counterfactual Reasoning \\in Autonomous Driving}


% The \author macro works with any number of authors. There are two commands
% used to separate the names and addresses of multiple authors: \And and \AND.
%
% Using \And between authors leaves it to LaTeX to determine where to break the
% lines. Using \AND forces a line break at that point. So, if LaTeX puts 3 of 4
% authors names on the first line, and the last on the second line, try using
% \AND instead of \And before the third author name.


\author{%
  Hayeon Oh and Donghwan Lee \\
  Department of Electrical Engineering \\
  Korea Advanced Institute of Science and Technology \\
  \texttt{\{hahao5252, donghwan\}@kaist.ac.kr} \\
}


\begin{document}


\maketitle


\begin{abstract}
End-to-end autonomous driving systems have been shown to suffer from shortcut learning and causal confusion, limiting their ability to generalize robustly under distribution shift. To mitigate these limitations, recent works based on Vision-Language Models (VLMs) have explored incorporating linguistic reasoning capabilities into driving decision-making, reporting improvements in driving performance. In particular, counterfactual (CF) reasoning has been proposed to enable self-reflective, high-level causal understanding. However, it remains unclear whether the language-based reasoning mitigates shortcut learning or merely shifts it from visual representations to linguistic patterns, creating an illusion of causal understanding. Moreover, existing approaches often conflate intervention-based reasoning with CF reasoning, leaving the effects of genuine CF reasoning unexplored. In this work, we examine what CF reasoning actually induces in VLM-based autonomous driving and observe that it shifts the action distribution toward more conservative behaviors while increasing uncertainty in reasoning and decision-making. This effect persists under reinforcement fine-tuning, suggesting that it is not merely a transient artifact. Furthermore, visual attention analysis shows that CF reasoning does not alter visual grounding, while causal intervention experiments reveal that masking causally relevant and non-causal factors in accident scenarios results in comparable decision change rates. Overall, these results suggest that CF reasoning does not induce genuine causal reasoning grounded in visual evidence. Instead, it acts as a linguistic shortcut that reshapes the action distribution, teaching models what to say but not how to reason, highlighting a gap between language-based reasoning and causal understanding in autonomous driving.
\end{abstract}

% TODO : 전체 연구에 대한 간략 요약 + 해당 연구의 당위성  
%% 연구 배경 : 왜 VLM, CoT, CoC, CF reasoning 연구가 등장하기 시작했는지 
%% question 제시, question 에 답하기 위해서 선정한 데이터셋과 metric, 그 이유
%% analysis 유형 논문임을 밝혀야 함
%% analysis 를 통해서 무얼 발견 했는지 (CF reasoning 이 self-reflection 을 유발하는 것처럼 보이지만, 언어적 차원에서 distribution shift 를 유발할 뿐, 그것이 마치 self-reflection 처럼 보이는 것. 그리고 이는 단순한 artifact 가 아니라 visual evidence 와 독립적이다.) 
\section{Introduction}

Autonomous driving based on the conventional end-to-end paradigm has been consistently shown to suffer from inherent limitations in causal reasoning due to issues such as shortcut learning and causal confusion. To address these challenges, recent studies have explored Vision-Language Models (VLMs), where language-based intermediate reasoning steps such as chain-of-thought (CoT) \cite{e2e-vlm:drivelm, e2e-vlm:drivegpt4} and chain-of-causation (CoC) \cite{causal:e2e-vlm:coc-vla, causal:e2e-vlm:alpamayo} have been reported to improve driving performance. In particular, counterfactual (CF) reasoning \cite{causal:e2e-vlm:omnidrive, causal:cf_vla, robotics:cast, robotics:chop} has been regarded as a promising approach for enabling higher-level causal understanding through self-reflective reasoning, thereby strengthening the expectation that VLMs can acquire genuine causal reasoning capabilities.

However, according to the formal definition in \cite{causal:ladder_of_causation}, CF reasoning defined in prior works \cite{causal:e2e-vlm:omnidrive, causal:cf_vla} often does not correspond to retrospective reasoning that reconstructs past events and reasons about alternative causes. Instead, it is more commonly limited to intervention-level reasoning that considers how actions might change under the current observation \cite{causal:e2e-vlm:simlingo}. In other words, rather than reasoning over unobserved alternative worlds, it is closer to action adjustment conditioned on the current input. Due to this discrepancy, it remains unclear whether the reported performance gains truly originate from causal reasoning.

This raises the fundamental question: what effects does genuine retrospective CF reasoning induce in VLMs for autonomous driving \textbf{(RQ1)}? More specifically, does it induce genuine causal reasoning, or does the introduction of language as an additional modality instead act as a shortcut that reshapes the action distribution \textbf{(RQ2)}?

% testbed dataset 이 accident dataset 인 이유
In this study, we address these questions by analyzing how CF reasoning, as an intermediate step prior to decision-making, influences model behavior and whether it leads to genuine causal understanding. Accordingly, evaluating these aspects requires scenarios with explicit causal factors. Furthermore, single-frame settings are insufficient to capture key causal cues, such as motion dynamics, trajectory evolution, and temporal risk accumulation, making temporal context essential. Accident video datasets provide a suitable setting for this evaluation, as they contain identifiable causal relationships and temporally ordered event progressions. Moreover, they enable retrospective counterfactual queries, such as "if the underlying cause had not been present, the accident would not have occurred," making them a suitable testbed for assessing genuine CF reasoning. 

% testbed dataset 이 meta-action 기반인 이유 (왜 meta-action 이 action distribution 변화, PPL 그리고 policy entropy 관측에 더 용이한지)
Beyond scenario design, an appropriate action space and evaluation metrics are required to examine how CF reasoning influences decision-making. Conventional trajectory-based evaluations primarily focus on driving performance, making them less suitable for our objective of analyzing changes in the model's decision-making behavior and its underlying action distribution. In particular, a trajectory lies in a continuous action space, whereas a meta-action represents a discrete driving behavior. Moreover, trajectory is less semantically interpretable than meta-actions, making it difficult to interpret changes in the likelihood of semantically equivalent decisions. Consequently, meta-action provides a more interpretable basis for analyzing changes in reasoning exploration. Therefore, we adopt meta-action as the decision space for analysis, following prior work\cite{e2e-vlm:senna, e2e-vlm:alphadrive, causal:cf_vla, e2e-vlm:driveagent-r1, e2e-vlm:senna-2}. This enables consistent evaluation across likelihood-based analysis using perplexity and direct examination of changes in distributional diversity through policy entropy. Accordingly, we construct a testbed dataset by integrating existing accident video datasets \cite{dataset:ccd, dataset:nexar, dataset:coool} with explicit causal factors, temporal context, and meta-action annotations. Further details of the dataset are provided in Appendix~\ref{appendix:dataset}.

% abstract 에서 모호하게 서술한 부분(결국 논문에서 발견한게 무엇인지)을 구체화
In conclusion, we reveal three key findings. First, we observe that CF reasoning induces a consistent shift toward more conservative action distributions, accompanied by increased entropy in language-based decision outputs, indicating that the model generates more diverse linguistic explanations while its decisions become biased toward conservative directions.The action distribution shift persists even after reinforcement learning fine-tuning (RFT), suggesting that it is not merely a training artifact. Second, we examine changes in visual attention with and without CF reasoning and find no meaningful differences in attention patterns, suggesting that the reasoning process operates independently of visual evidence. Third, in causal intervention experiments, VLMs exhibit similar levels of decision change regardless of whether causal or non-causal objects are masked, unlike humans selectively change their decisions only when causal factors are masked. This indicates that CF reasoning does not enable the model to selectively identify causal factors.

% 추가 분석과 더불어서 본 논문에서의 findings 를 요약
Overall, these results indicate that CF reasoning primarily reshapes action distributions rather than improving the grounding of causal reasoning in visual evidence. We further examine whether this effect extends to general driving scenarios and low-level decisions, as in prior work reporting performance gains from CF reasoning\cite{causal:e2e-vlm:omnidrive}. An ablation on CF reasoning data reveals that models trained with CF reasoning consistently exhibit lower rates of lateral maneuvers than those trained without it, suggesting that this shift also manifests under existing CF reasoning paradigms. Given the relatively low prevalence of lateral maneuvers in the evaluation trajectories, this shift may align with the dataset's action distribution, potentially explaining the reported performance gains. Taken together, our findings indicate that CF reasoning may not induce causal reasoning capabilities, but rather functions as a linguistic shortcut that shifts the existing policy distribution.


\section{Preliminaries}
\label{sec:preliminaries}
% TODO : 기존 자율주행에서의 metric 의 한계 + perpelxity 와 pass@k 를 쓰는 이유
% \subsection{The Misalignment between Curve Fitting and Reasoning.}
\subsection{Limitation of Driving Metrics for Evaluating Reasoning}

% 자율주행 연구에서는 모델의 성능을 expert trajectory와의 유사도 \cite{dataset:nusc}나 주행 성공률과 같은 driving performance 지표를 통해 평가하는 것이 일반적이다 \cite{dataset:navsim, dataset:bench2drive}. 이러한 지표들은 주행 행동을 정량적으로 비교하는 데에는 유용하지만, 모델이 실제로 어떤 추론 과정을 거쳐 결정을 내렸는지를 직접적으로 드러내지는 못한다. 특히 최근 연구들에서는 높은 주행 성능이 반드시 시각적 이해나 인과적 추론 능력의 향상으로부터 비롯되지 않을 수 있음을 보여준다. 예를 들어, ego status나 navigation information과 같은 비시각적 입력만으로도 상당한 수준의 주행 성능을 달성할 수 있으며 \cite{causal:e2e:ad-mlp, e2e-ad:is-ego-status}, 이는 VLM 기반 연구에서도 관찰된다\cite{causal:e2e-vlm:omnidrive, causal:e2e-vlm:more_than_meets_the_eye}. 이러한 결과는 모델이 환경에 대한 인과적 이해를 수행하기보다는 데이터에 내재된 상관관계나 shortcut 패턴에 의존하여 의사결정을 수행하고 있을 가능성을 시사한다. 다시 말해, 동일한 주행 성능의 향상이 관찰되더라도 그것이 인과적 추론 능력의 향상에 기인한 것인지, 혹은 데이터 편향이나 shortcut learning에 의한 것인지 구분하기 어렵다는 confounding 문제가 존재한다. 따라서 기존의 주행 성능 기반 평가만으로는 CF reasoning의 효과를 신뢰성 있게 분석하는 데 근본적인 한계가 있다. 이를 해결하기 위해 본 연구에서는 잠재적인 추론 특성을 정량적으로 평가할 수 있는 지표들을 채택한다. Detailed definitions for each metric are provided in Appendix \ref{appendix:implementations}.
Metrics such as displacement error between predicted and expert trajectories \cite{dataset:nusc} and overall driving success rates \cite{dataset:navsim, dataset:bench2drive} are widely used for evaluating autonomous driving systems. However, while these metrics are effective for quantitative performance comparison, they do not directly reflect the underlying reasoning process behind model decisions. Recent studies further show that strong driving performance does not necessarily imply improved visual understanding or causal reasoning ability. In particular, models can achieve competitive results using only non-visual inputs such as ego status and navigation signals \cite{causal:e2e:ad-mlp, e2e-ad:is-ego-status}, and similar behaviors have also been observed in VLM-based approaches \cite{causal:e2e-vlm:omnidrive, causal:e2e-vlm:more_than_meets_the_eye}. These findings suggest that current performance gains may stem from spurious correlations or shortcut patterns in the data rather than genuine causal reasoning, making it difficult to isolate the effect of CF reasoning using standard driving metrics alone. Consequently, it remains unclear whether performance gains reflect improved reasoning or dataset-induced shortcuts, making it difficult to isolate the effect of CF reasoning. To address this issue, we adopt evaluation metrics designed to quantify latent reasoning behavior. Detailed definitions are provided in Appendix~\ref{appendix:implementations}.

% TODO : 세 가지 METRICS 가 무엇인지, 왜 쓰는지 피력. 
% 이런 metric 을 쓴다는 주장이 자율주행 연구자들에게는 생소할 수도 있음. 따라서 참고 문헌도 배치할 것 
\subsection{Metrics for Reasoning Capacity and Exploration}


% TODO : pass@k 는 action distribution shift 보여주는 용도 
\paragraph{Pass@k.} 
% 이것은 상위 k개의 생성 결과 내에 정답이 포함되는지를 측정하며, 일반적으로 모델이 특정 문제에 대해 올바른 답을 생성할 수 있는 잠재적 능력을 평가한다. 기존의 precision, recall과 같은 분류 기반 지표는 단일 예측 결과에 대한 정답 여부만을 평가하기 때문에, 모델이 실제로 어떤 선택지를 고려했는지, 그리고 그 중 어떤 결정을 선택했는지를 구분하지 못한다. 반면, Pass@k는 하나의 입력에 대해 생성된 다수의 후보를 기반으로 평가를 수행함으로써, 모델이 생성 가능한 의사결정의 범위와 실제 선택 간의 차이를 분리하여 분석할 수 있게 한다. 특히, Pass@1은 모델의 최종 의사결정을 반영하는 반면, k가 커질수록 모델이 생성할 수 있는 잠재적인 의사결정 집합을 근사한다. 이러한 구분은 CF reasoning이 새로운 선택지를 생성하는지, 혹은 기존 선택지들 간의 선택 확률을 재조정하는지를 구분하여 분석할 수 있게 한다.
It measures whether a correct answer is included within the top-k generated outputs, and is commonly used to evaluate a model’s potential capability to solve a given problem \cite{metric:passk, limit-rlvr:does, metric:passk:example:atlas}. Unlike single-prediction metrics such as precision and recall, Pass@k evaluates multiple candidates per input, allowing us to distinguish between generation capability and selection behavior. In particular, Pass@1 reflects the model’s final decision, while larger values of k approximate the set of feasible decisions that the model can generate. This allows us to determine whether CF reasoning introduces new decisions or instead reshapes the distribution of existing actions.


% TODO : perplexity 가 무엇인지, 왜 쓰이는지, 우리는 어떻게 쓰는지 
%% In this work, we compute $\text{PPL}_{base}(Y_{\text{\text{GT}}}|x)$ and $\text{PPL}_{CF}(Y_{\text{\text{GT}}}|x)$.
\paragraph{Perplexity.} 
% 만약 CF reasoning에 따라 의사결정 결과의 분포가 변화한다면, 이러한 변화가 단순한 확률 분포의 재조정에 불과한 것인지, 혹은 의사결정 생성 과정에 대한 실질적인 개입에 기인한 것인지를 구분할 필요가 있다. 이를 위해 본 연구에서는 주어진 출력 시퀀스에 대해 모델이 부여하는 likelihood를 측정하는 perplexity를 추가적으로 분석한다. 
If the distribution of decision outcomes shifts under CF reasoning, it is necessary to determine whether such changes merely reflect a redistribution of probabilities or arise from substantive intervention in the decision generation process. To this end, we additionally analyze perplexity, which measures the likelihood assigned by a model to a given output sequence, and is defined in Appendix \ref{lst:perplexity}. This quantifies how naturally a model generates a given output sequence under its conditional distribution. This aims to distinguish whether the observed action distribution shift is simply a change in output probabilities or a consequence of CF reasoning acting as a conditioning factor in decision making. Specifically, if CF reasoning is merely an auxiliary process unrelated to decision making, both models are expected to assign similar likelihoods to the same Y. In contrast, if CF reasoning substantively contributes to decision generation, the model generating response with CF reasoning may assign a lower likelihood to the same output when evaluated without conditioning on the reasoning process, which would be reflected as higher perplexity.


% TODO : GRPO, GDPO 를 통해, CF reasoning 의 효과가 artifact 가 아님을 보이려고 함을 피력
\paragraph{RFT as a Verification Framework.}
The pass@k and perplexity analyses presented above are conducted on supervised fine-tuned (SFT) models that are instruction-tuned to perform CF reasoning prior to decision making. However, it remains unclear whether the observed behavioral patterns are merely artifacts. To address this question, we employ reinforcement fine-tuning (RFT) as a diagnostic framework. In autonomous driving, RFT enables policy optimization through explicitly designed reward functions that align learning objectives with driving performance metrics \cite{e2e-vlm:autovla, e2e-vlm:alphadrive, e2e-vlm:autodrive-r2, causal:e2e-vlm:alpamayo}. In this work, we adopt GRPO \cite{rlvr:grpo} and its variant GDPO \cite{rlvr:gdpo}, both of which generate $G$ outputs per query and update the model based on relative performance within each group. Their shared objective function is defined as follows:
$$\mathcal{J}(\theta) = \mathbb{E}_{q, \{o_i\}} \left[ \frac{1}{G} \sum_{i=1}^{G} \min \left( w_i A_i, \text{clip}(w_i, 1-\epsilon, 1+\epsilon) A_i \right) - \beta \mathbb{D}_{\mathrm{KL}}(\pi_\theta \| \pi_{\text{ref}}) \right]$$
where $w_i$ denotes the importance sampling ratio and $A_i$ represents the group-relative advantage as defined in Appendix \ref{appendix:implementations}. Importantly, the KL divergence constraint prevents the updated policy $\pi_\theta$ from deviating excessively from the reference policy $\pi_{\text{ref}}$, thereby stabilizing the optimization process and preserving the fundamental characteristics of the original policy \cite{limit-rlvr:does, limit-rlvr:echo_chamber, limit-rlvr:rethinking-reflection}. Thus, RFT can be a surrogate for probing the latent behavioral properties of the SFT reference model. In particular, if the behavioral tendencies induced by CF reasoning persist throughout RFT optimization, this suggests that the observed effects are unlikely to be incidental artifacts of the data distribution, but rather reflect more structured learning signals.

\paragraph{Policy Entropy as a Real-Time Exploration Probe.}
correct English
Policy entropy, defined in Appendix \ref{eq:policy_entropy} and recorded throughout RFT, measures token-level decision diversity during generation and has recently been adopted as a real-time indicator of exploration capacity. \cite{policy-entropy:cui2025entropy, policy-entropy:liu2025prorl}. Higher entropy indicates that the model explores a broader decision space, whereas a sharp decrease in entropy suggests convergence toward specific action patterns, i.e., reduced exploration. In this work, we leverage policy entropy as an additional diagnostic signal for a more structural analysis of the effect of CF reasoning. Specifically, we aim to examine whether CF reasoning structurally alters the model’s exploratory behavior, or whether the previously observed effects are merely incidental artifacts of the SFT training setup. If CF reasoning indeed acts as an inductive bias that expands the model’s decision-making exploration, it should consistently manifest as higher entropy in CF models, in alignment with the observed action distribution shifts and perplexity changes \cite{policy-entropy:reasoning_w_exploration}.
% CF reasoning을 포함하는 모델은 구조적으로 더 긴 completion을 생성하므로, 길이 차이를 통제하지 않으면 단순한 시퀀스 길이의 차이가 entropy 비교를 왜곡할 수 있다. 본 연구에서는 이러한 confound 를 억제하기 위해 length-normalized entropy를 통해 이러한 길이 편향을 제거하고, CF reasoning 이 RFT 학습 이후에도 전반의 탐색 분포에 미치는 지를 관찰한다.

% TODO : (앞선 문단에서 perplexity 와 pass@k 를 쓰는 이유와 연결지어서) 평가 기준 시나리오가 왜 사고 시나리오인지, decision 을 low-level trajectory 가 아니라 high-level decision 인 meta-action 을 기준 삼아서 평가하는지, CF reasoning 의 효과를 평가하기 위해 instruction fine-tuning 을 했다는 점을 명시
\section{Problem Formulation and Experimental Setup}
We define each sample in the dataset as $Y = (a, e, d, r, y)$, where $a \in \{0,1\}$ denotes accident occurrence, $e \in \{0,1\}$ indicates ego-vehicle involvement, $d$ represents the scene description extracted from the video, $r$ denotes the CF reasoning component, and $y$ corresponds to the final meta-action decision. 
Furthermore, we define three longitudinal meta-actions: \texttt{Stop}, \texttt{Decelerate}, and \texttt{Keep}. These span the full range of longitudinal responses from immediate braking to speed maintenance and are mutually exclusive and collectively exhaustive with respect to longitudinal behavior in accident scenarios, making them appropriate for analyzing how CF reasoning affects the action distribution. To analyze the effect of CF reasoning, we fine-tune VLMs to generate responses in the following structured format: accident occurrence, ego involvement, scene descriptions, CF reasoning, and final decision. This ordering reflects the causal structure of retrospective reasoning: the model first reconstructs what happened and why, then determines the appropriate response, encouraging it to condition its final action on CF reasoning. As a baseline, we fine-tune a model with the same structure but without the CF reasoning component, enabling a controlled comparison of the effects of CF reasoning on decision making. 



\section{Analysis of Counterfactual Reasoning}
\label{sec:analysis}

% TODO : main claim 이자 분석을 소개하는 부분. 
% pass@1, pass@8 을 통해 action distribution shift 라는 현상을 주장
% perplexity 를 통해 overconfidence 라는 internal state 를 관측 
\subsection{Counterfactual Reasoning Induces Behavioral Shifts}
\label{subsec:sft_analysis}

% ===============================
% Table 1: SFT Stage Results
% ===============================

\begin{table*}[h]
\centering
\small
\resizebox{\textwidth}{!}{
\begin{tabular}{l c | cc cc | ccc}
\toprule
\multirow{2}{*}{\textbf{Model}} & \multirow{2}{*}{\textbf{Metric}}
  & \multicolumn{2}{c}{\textbf{Accident Recog.}}
  & \multicolumn{2}{c|}{\textbf{Ego Inclusion}}
  & \multicolumn{3}{c}{\textbf{Decision}} \\
% \cmidrule{7-9}
 & & Yes & No & Involv. & Not In. & Stop & Decel. & Keep \\
\midrule

\multirow{6}{*}{\textbf{InternVL2-1B}}
 & P@1 w/o & 47.5 & 64.4 & 45.8 & 64.4 & 20.4 & 29.2 & 78.0 \\
 & P@1 w/  & 60.3 & 50.6 & 60.3 & 50.6 & 26.2 & 24.9 & 67.2 \\
 & $\Delta$
   & \cellcolor{blue!15}+12.8 & \cellcolor{red!15}-13.8
   & \cellcolor{blue!15}+14.5 & \cellcolor{red!15}-13.8
   & \cellcolor{blue!15}+5.8 & \cellcolor{red!15}-4.3
   & \cellcolor{red!15}-10.8 \\
\cmidrule{2-9}
 & P@16 w/o & 99.4 & 99.7 & 99.4 & 99.7 & 87.2 & 90.3 & 99.8 \\
 & P@16 w/  & 99.9 & 99.9 & 99.9 & 99.9 & 93.3 & 93.6 & 100.0 \\
 & $\Delta$
   & \cellcolor{red!15}+0.5  & \cellcolor{blue!15}+0.2
   & \cellcolor{red!15}+0.5  & \cellcolor{blue!15}+0.2
   & \cellcolor{blue!15}+6.1 & \cellcolor{blue!15}+3.3
   & \cellcolor{blue!15}+0.2 \\
\midrule

\multirow{6}{*}{\textbf{InternVL2-2B}}
 & P@1 w/o & 38.8 & 74.1 & 38.8 & 74.1 & 17.3 & 22.8 & 83.2 \\
 & P@1 w/  & 54.3 & 66.1 & 54.3 & 66.1 & 26.5 & 26.2 & 78.8 \\
 & $\Delta$
   & \cellcolor{blue!15}+15.5 & \cellcolor{red!15}-8.0
   & \cellcolor{blue!15}+15.5 & \cellcolor{red!15}-8.0
   & \cellcolor{blue!15}+9.2 & \cellcolor{blue!15}+3.4
   & \cellcolor{red!15}-4.4 \\
\cmidrule{2-9}
 & P@16 w/o & 92.9 & 99.9 & 92.9 & 99.9 & 76.0 & 80.3 & 100.0 \\
 & P@16 w/  & 95.8 & 99.0 & 95.8 & 99.0 & 78.0 & 84.1 & 99.7 \\
 & $\Delta$
   & \cellcolor{blue!15}+2.9  & \cellcolor{red!15}-0.9
   & \cellcolor{blue!15}+2.9  & \cellcolor{red!15}-0.9
   & \cellcolor{blue!15}+2.0 & \cellcolor{blue!15}+3.8
   & \cellcolor{red!15}-0.3 \\
\midrule

\multirow{6}{*}{\textbf{InternVL2.5-1B}}
 & P@1 w/o & 53.1 & 65.9 & 53.1 & 65.9 & 20.1 & 25.4 & 74.3 \\
 & P@1 w/  & 56.1 & 63.2 & 56.1 & 63.2 & 27.8 & 27.4 & 73.8 \\
 & $\Delta$
   & \cellcolor{blue!15}+3.0  & \cellcolor{red!15}-2.7
   & \cellcolor{blue!15}+3.0  & \cellcolor{red!15}-2.7
   & \cellcolor{blue!15}+7.7 & \cellcolor{blue!15}+2.0
   & \cellcolor{red!15}-0.5 \\
\cmidrule{2-9}
 & P@16 w/o & 99.1 & 99.4 & 99.1 & 99.4 & 85.0 & 87.4 & 99.8 \\
 & P@16 w/  & 99.1 & 98.8 & 99.1 & 98.8 & 87.2 & 88.7 & 99.4 \\
 & $\Delta$
   & \cellcolor{blue!15}+0.0  & \cellcolor{red!15}-0.6
   & \cellcolor{blue!15}+0.0  & \cellcolor{red!15}-0.6
   & \cellcolor{blue!15}+2.2 & \cellcolor{blue!15}+1.3
   & \cellcolor{red!15}-0.4 \\
\midrule

\multirow{6}{*}{\textbf{Qwen2.5-VL-3B}}
 & P@1 w/o & 60.9 & 49.4 & 54.3 & 52.7 & 21.7 & 38.2 & 59.3 \\
 & P@1 w/  & 75.6 & 32.6 & 74.6 & 33.0 & 45.7 & 36.7 & 45.1 \\
 & $\Delta$
   & \cellcolor{blue!15}+14.7 & \cellcolor{red!15}-16.8
   & \cellcolor{blue!15}+20.3 & \cellcolor{red!15}-19.7
   & \cellcolor{blue!15}+24.0 & \cellcolor{red!15}-1.5
   & \cellcolor{red!15}-14.2 \\
\cmidrule{2-9}
 & P@16 w/o & 98.5 & 90.3 & 98.1 & 93.3 & 76.7 & 94.9 & 97.5 \\
 & P@16 w/  & 100.0 & 86.0 & 100.0 & 86.5 & 97.1 & 94.6 & 92.5 \\
 & $\Delta$
   & \cellcolor{blue!15}+1.5  & \cellcolor{red!15}-4.3
   & \cellcolor{blue!15}+1.9  & \cellcolor{red!15}-6.8
   & \cellcolor{blue!15}+20.4 & \cellcolor{red!15}-0.3
   & \cellcolor{red!15}-5.0 \\
\midrule

\multirow{6}{*}{\textbf{VideoLLaMA2-7B}}
 & P@1 w/o & 47.9 & 53.3 & 47.9 & 53.3 & 16.9 & 26.7 & 69.7 \\
 & P@1 w/  & 72.3 & 28.7 & 71.7 & 30.1 & 30.0 & 30.0 & 53.8 \\
 & $\Delta$
   & \cellcolor{blue!15}+24.4 & \cellcolor{red!15}-24.6
   & \cellcolor{blue!15}+23.8 & \cellcolor{red!15}-23.2
   & \cellcolor{blue!15}+13.1 & \cellcolor{blue!15}+3.3
   & \cellcolor{red!15}-15.9 \\
\cmidrule{2-9}
 & P@16 w/o & 99.3 & 100.0 & 99.3 & 100.0 & 74.8 & 95.4 & 100.0 \\
 & P@16 w/  & 100.0 & 91.8 & 100.0 & 92.9 & 86.6 & 97.2 & 98.6 \\
 & $\Delta$
   & \cellcolor{blue!15}+0.7  & \cellcolor{red!15}-8.2
   & \cellcolor{blue!15}+0.7  & \cellcolor{red!15}-7.1
   & \cellcolor{blue!15}+11.8 & \cellcolor{blue!15}+1.8
   & \cellcolor{red!15}-1.4 \\
\bottomrule
\end{tabular}
}
\caption{
  Pass@1 and Pass@16 comparison w/ and w/o counterfactual reasoning in SFT.
  \textbf{Conservative} decisions (Stop / Decel.) are grouped under a shared header.
  In the $\Delta$ rows, \colorbox{blue!15}{\textbf{bold blue}} highlights gain
  and \colorbox{red!15}{\textbf{bold red}} highlights reduction,
  together indicating a consistent shift toward conservative decision-making
  when counterfactual reasoning is included.
}
\label{tab:passk_sft}
\end{table*}


In this section, we evaluate the impact of CF reasoning on VLMs during the supervised fine-tuning (SFT). To verify whether the observed findings consistently generalize across different VLMs, we extend our evaluation to a diverse set of models, including InternVL2-1B, InternVL2-2B, InternVL2.5-1B, Qwen2.5-VL-3B, and VideoLLaMA2-7B, with detailed selection criteria provided in Appendix \ref{appendix:model_selection_criteria}.

% TODO : CF reasoning 이 action distribution shift 를 유발한다는 주장의 개요 
% TODO : decision 에만 집중해서 설명하지 말 것. 즉, recognition 과 inclusion 결과 먼저 언급 필요 
\paragraph{Action Distribution Shift.} Pass@1 reflects the model’s single-shot decision behavior, capturing how the most confident action changes depending on the presence of CF reasoning. Across most models, CF reasoning consistently shifts predictions toward more conservative actions such as \texttt{Stop} and \texttt{Decelerate}, while reducing the proportion of \texttt{Keep}, suggesting a directional bias in decision preference rather than a simple increase in output variability. Extending this observation, Pass@16 approximates the broader set of reachable decisions under multi-sampling and similarly exhibits a shift toward conservative actions. This indicates that the effect is not limited to stochastic single-generation behavior, but instead reflects a systematic reorganization of the model’s underlying action distribution. Overall, these results suggest that CF reasoning induces a structurally embedded and persistent change in the model’s decision space, which becomes more pronounced in larger models, implying stronger internalization of the induced reasoning structure.

% TODO : perplexity 는 reasoning capacity 의 확장을 표현하는게 아님 
% pass@k 의 결과의 원인을 설명하는 단계 
\paragraph{Does CF Reasoning Intervene in Decision Generation?} 

However, it remains unclear whether the observed distribution shift is a surface-level artifact caused by the co-occurrence of CF reasoning tokens or whether it reflects a structural intervention of CF reasoning in the decision generation process that reshapes the model’s action distribution. To investigate this, we evaluate both models on a shared output, namely the common text $Y^{(-r)}_{\text{GT}}$, obtained by excluding the CF reasoning component, and compare their perplexity. If CF reasoning merely serves as an auxiliary output unrelated to decision generation, both models should assign similar likelihoods to $Y^{(-r)}_{\text{GT}}$. In contrast, if CF reasoning functions as a conditioning variable that actively influences decision generation, the CF model is expected to assign lower likelihood due to the removal of a key conditioning component, resulting in higher perplexity. Accordingly, the gap between $\text{PPL}_{\text{base}}(Y^{(-r)}_{\text{GT}})$ and $\text{PPL}_{\text{CF}}(Y^{(-r)}_{\text{GT}})$ serves as a diagnostic signal for structural intervention. As shown in Figure \ref{fig:ppl_sft}, the CF model consistently exhibits higher perplexity across all action classes, suggesting that CF reasoning acts as a conditioning factor that reshapes the decision-generation distribution rather than serving as a purely auxiliary output.


\begin{figure}[h]
    \centering
    \includegraphics[width=\textwidth]{figure/ppl/ppl_classwise_sft.pdf}
    \caption{Class-wise perplexity comparison between the models w/o \textbf{(blue)} and w/ CF \textbf{(orange)} reasoning on the shared output $Y^{(-r)}_{\text{GT}}$, corresponding to $\text{PPL}_{\text{base}}(Y^{(-r)}_{\text{GT}})$ and $\text{PPL}_{\text{CF}}(Y^{(-r)}_{\text{GT}})$, respectively.}
    % \label{fig:ppl_classwise_sft}
    \label{fig:ppl_sft}
\end{figure}



% TODO : 왜 강화학습 파인튜닝을 했는지 서술하면서 이전 문단과 흐름이 자연스러워야함
\subsection{Persistence of Behavioral Shifts under Reinforcement Fine-Tuning}
\label{subsec:rl_analysis}

However, it remains unclear whether these changes are merely artifacts or reflect a structural bias that persists throughout training. If CF reasoning serves only as a superficial signal for correct answers, its effect should diminish under reinforcement learning as reward optimization proceeds. In contrast, if it fundamentally reshapes the policy space, the induced bias should persist after reinforcement learning. To investigate this, we conduct RFT using the SFT models as reference models. We evaluate InternVL2.5-1B, which exhibits the smallest bias based on the combined change in \texttt{stop} and \texttt{decel.} under Pass@16, and Qwen2.5-VL-3B, which shows the largest bias. Details of the reward functions are provided in Appendix \ref{appendix:reward_functions}.
% ===============================
% Table 2: RFT Stage Results
% ===============================
\begin{table*}[h]
\centering
\small
\resizebox{\textwidth}{!}{
\begin{tabular}{l c | cc cc | ccc}
\toprule
\multirow{2}{*}{\textbf{Model}} & \multirow{2}{*}{\textbf{Metric}}
  & \multicolumn{2}{c}{\textbf{Accident Recog.}}
  & \multicolumn{2}{c|}{\textbf{Ego Inclusion}}
  & \multicolumn{3}{c}{\textbf{Decision}} \\
% \cmidrule{7-9}
 & & Yes & No & Involv. & Not In. & Stop & Decel. & Keep \\
\midrule

\multirow{6}{*}{\textbf{InternVL2.5-1B (GRPO)}}
 & P@1 w/o & 36.6 & 79.9 & 36.0 & 80.4 & 16.9 & 16.7 & 87.2 \\
 & P@1 w/  & 39.6 & 73.5 & 39.6 & 73.5 & 24.6 & 17.9 & 80.3 \\
 & $\Delta$
   & \cellcolor{blue!15}+3.0  & \cellcolor{red!15}-6.4
   & \cellcolor{blue!15}+3.6  & \cellcolor{red!15}-6.9
   & \cellcolor{blue!15}+7.7  & \cellcolor{blue!15}+1.2
   & \cellcolor{red!15}-6.9 \\
\cmidrule{2-9}
 & P@16 w/o & 92.9 & 99.7 & 92.9 & 99.7 & 73.8 & 75.1 & 100.0 \\
 & P@16 w/  & 96.3 & 99.9 & 96.3 & 99.9 & 82.7 & 84.4 &  99.8 \\
 & $\Delta$
   & \cellcolor{blue!15}+3.4  & \cellcolor{blue!15}+0.2
   & \cellcolor{blue!15}+3.4  & \cellcolor{blue!15}+0.2
   & \cellcolor{blue!15}+8.9  & \cellcolor{blue!15}+9.3
   & \cellcolor{red!15}-0.2 \\
\midrule

\multirow{6}{*}{\textbf{Qwen2.5-VL-3B (GRPO)}}
 & P@1 w/o & 49.9 & 56.8 & 47.5 & 58.5 & 17.9 & 35.9 & 65.8 \\
 & P@1 w/  & 71.0 & 32.6 & 67.3 & 35.3 & 41.9 & 32.3 & 42.4 \\
 & $\Delta$
   & \cellcolor{blue!15}+21.1 & \cellcolor{red!15}-24.2
   & \cellcolor{blue!15}+19.8 & \cellcolor{red!15}-23.2
   & \cellcolor{blue!15}+24.0 & \cellcolor{red!15}-3.6
   & \cellcolor{red!15}-23.4 \\
\cmidrule{2-9}
 & P@16 w/o & 95.7 & 94.0 & 95.7 & 94.8 & 69.6 & 92.3 & 96.7 \\
 & P@16 w/  & 100.0 & 89.0 & 100.0 & 92.9 & 94.6 & 92.6 & 92.8 \\
 & $\Delta$
   & \cellcolor{blue!15}+4.3  & \cellcolor{red!15}-5.0
   & \cellcolor{blue!15}+4.3  & \cellcolor{red!15}-1.9
   & \cellcolor{blue!15}+25.0 & \cellcolor{blue!15}+0.3
   & \cellcolor{red!15}-3.9 \\
\midrule

\multirow{6}{*}{\textbf{Qwen2.5-VL-3B (GDPO)}}
 & P@1 w/o & 45.4 & 62.5 & 42.7 & 64.9 & 13.4 & 29.0 & 70.5 \\
 & P@1 w/  & 74.7 & 38.2 & 68.8 & 41.7 & 44.4 & 68.8 & 44.4 \\
 & $\Delta$
   & \cellcolor{blue!15}+29.3 & \cellcolor{red!15}-24.3
   & \cellcolor{blue!15}+26.1 & \cellcolor{red!15}-23.2
   & \cellcolor{blue!15}+31.0 & \cellcolor{blue!15}+39.8
   & \cellcolor{red!15}-26.1 \\
\cmidrule{2-9}
 & P@16 w/o & 95.2 & 97.5 & 95.2 & 98.2 & 68.7 & 91.0 & 98.9 \\
 & P@16 w/  & 100.0 & 87.5 & 100.0 & 94.3 & 93.9 & 93.6 & 89.4 \\
 & $\Delta$
   & \cellcolor{blue!15}+4.8  & \cellcolor{red!15}-10.0
   & \cellcolor{blue!15}+4.8  & \cellcolor{red!15}-3.9
   & \cellcolor{blue!15}+25.2 & \cellcolor{blue!15}+2.6
   & \cellcolor{red!15}-9.5 \\
\bottomrule
\end{tabular}
}
\caption{
  Pass@1 and Pass@16 comparison w/ and w/o counterfactual reasoning after RFT (GRPO / GDPO).
  \textbf{Conservative} decisions (Stop / Decel.) are grouped under a shared header.
  In the $\Delta$ rows, \colorbox{blue!15}{\textbf{bold blue}} highlights gain
  and \colorbox{red!15}{\textbf{bold red}} highlights reduction.
  The consistent shift toward conservative decision-making across both GRPO and GDPO
  confirms that the observed phenomenon is not a training artifact.
}
\label{tab:passk_rft}
\end{table*}

\paragraph{Persistence of CF-Induced Behavioral Shift under Reinforcement Learning.} To examine whether the distribution shift could be mitigated through reinforcement learning, we conducted GRPO on InternVL2.5-1B and Qwen2.5-VL-3B, taking into account both algorithmic and model generality. Subsequently, we conducted GDPO on Qwen2.5-VL-3B, which exhibited the strongest bias during SFT, to further validate algorithmic generality. As shown in Table~\ref{tab:passk_rft}, the decision distribution changes induced by CF reasoning were maintained after RFT. If the patterns observed during SFT were mere training artifacts, they should have been corrected during reinforcement learning through reward optimization; however, the results indicate otherwise. In other words, the policy bias formed by CF reasoning persists after RFT. See Appendix \ref{appendix:passk_sft} and \ref{appendix:passk_rft} for detail analysis.


\begin{figure*}[h]
\centering
\begin{minipage}{0.48\textwidth}
Figure \ref{fig:ppl_rft} shows class-wise perplexity after RFT. Across all action classes, $\text{PPL}_{\text{CF}}(Y^{(-r)}_{\text{GT}})$ remains consistently higher than $\text{PPL}_{\text{base}}(Y^{(-r)}_{\text{GT}})$, following the same trend observed in SFT. Taken together, both pass@k and perplexity results indicate that the distributional effects induced by CF reasoning persist after reinforcement learning, suggesting a persistent structural influence on the model’s generative distribution rather than a superficial behavioral change.
\end{minipage}
\hfill
\begin{minipage}{0.48\textwidth}
    \centering
    \includegraphics[width=\linewidth]{figure/ppl/ppl_classwise_rft.pdf}
    \caption{Class-wise perplexity comparison between the models w/o \textbf{(blue)} and w/ CF \textbf{(orange)} reasoning  on the shared output $Y^{(-r)}_{\text{GT}}$ after RFT.}
    \label{fig:ppl_rft}
\end{minipage}
\end{figure*}


\paragraph{Policy Entropy as a Dynamic Probe of Exploration.}
Pass@k and perplexity are metrics that evaluate the resulting action distribution after training convergence; however, they do not directly capture how the policy is formed during learning or how the exploration space evolves over time. To address this limitation, we log and analyze policy entropy throughout the entire RFT process. Since RFT incrementally updates the policy under a KL-divergence constraint with respect to the SFT reference model, it provides a suitable setting for dynamically tracking changes in the action distribution over the full training trajectory. Empirically, models incorporating CF reasoning consistently exhibit higher policy entropy throughout RFT. This indicates that CF reasoning alters the exploration structure of the policy during training. In particular, when considered alongside the shift toward conservative actions such as \texttt{Stop} and \texttt{Decelerate} observed in Pass@k, this increase in entropy does not reflect a mere rise in randomness. Instead, it can be interpreted as a redistribution of probability mass from a previously over-concentrated preference on \texttt{Keep} toward a more balanced allocation over multiple valid actions. Overall, the policy entropy analysis provides complementary evidence that CF reasoning affects not only the relative preference among actions but also the exploration structure of the policy during training, thereby reinforcing the findings from Pass@k and perplexity regarding its structural impact.

\begin{figure*}[h]
    \centering

    % ---------- Top row ----------
    \makebox[\textwidth][c]{%
        \begin{subfigure}{0.32\textwidth}
            \centering
            \includegraphics[width=\linewidth]{figure/entropy/policy_entropy_internvl2_5_1b.pdf}
            \caption{InternVL2.5-1B (GRPO)}
        \end{subfigure}
        \begin{subfigure}{0.32\textwidth}
            \centering
            \includegraphics[width=\linewidth]{figure/entropy/policy_entropy_qwen2_5_vl.pdf}
            \caption{Qwen2.5-VL-3B (GRPO)}
        \end{subfigure}
        \hspace{1em}
        \begin{subfigure}{0.32\textwidth}
            \centering
            \includegraphics[width=\linewidth]{figure/entropy/policy_entropy_qwen2_5_vl_gdpo.pdf}
            \caption{Qwen2.5-VL-3B (GDPO)}
        \end{subfigure}
    }


    \caption{Policy entropy over training steps during RFT.}
    \label{fig:entropy}
\end{figure*}



\subsection{CF Reasoning Decouples from Visual Grounding}
\label{subsec:grounding_analysis}

In summary, we indirectly examined the impact of CF reasoning on decision generation by focusing on the model’s generated text outputs and their probabilistic characteristics. However, the metrics used in the previous analysis are suitable for studying latent behavioral reasoning, but they do not directly evaluate alignment with causal structure. Therefore, we conduct a more direct investigation to determine whether the observed effects arise from genuine causal reasoning induced by CF reasoning or from linguistic bias in this section.


% TODO : 왜 TAM ? 왜 두 가지 모델만?
% TODO : 세 가지 지표가 의미하는 바 
\paragraph{Token Activation Map Analysis.} If CF reasoning couples with perceptual awareness, the model’s decisions should rely more strongly on the visual causes of accidents, with attention correspondingly focusing on relevant risk-related regions. To examine whether CF reasoning induces decision-making bias through genuine visual risk perception, we conduct a visual attribute analysis using Token Activation Maps (TAM) \cite{attention:tam}, which isolate token-level activations to estimate the visual regions on which each output token depends. Unlike methods affected by inter-token interactions, TAM provides a more stable measure of the visual cues used by specific reasoning tokens. Based on this, we compute three metrics: (i) token-level IoU ($IoU_{tok}$), which measures how precisely individual causal keyword tokens align with accident-causing objects; (ii) aggregated IoU ($IoU_{agg}$), which evaluates whether attention remains consistently focused on the target object over the entire reasoning process; and (iii) grounding F1 ($F1_{gnd}$), which jointly accounts for precision and recall to balance the completeness and accuracy of attention localization described in Appendix ~\ref{appendix:implementations_iou_metrics}. As shown in Table~\ref{tab:iou}, IoU values vary across models but remain largely unaffected by prompt variations, indicating consistent spatial grounding patterns. However, no clear alignment between attention and accident-causing objects is observed, suggesting that the CF-induced action distribution shifts may occur independently of visual evidence as shown in Figure \ref{fig:attention_map}.

\begin{table*}[h] 
\centering
\resizebox{\textwidth}{!}{%
\begin{tabular}{@{}lcccccccccccc@{}}
\toprule
\multirow{3}{*}{} & \multicolumn{3}{c}{\textbf{Ablation}} & \multicolumn{3}{c}{\textbf{perception}} & \multicolumn{3}{c}{\textbf{decision}} & \multicolumn{3}{c}{\textbf{counterfactual}} \\
\cmidrule(lr){2-4} \cmidrule(lr){5-7} \cmidrule(lr){8-10} \cmidrule(lr){11-13}
 & \textbf{Type} & \textbf{CF} & \textbf{RFT} & \textbf{$IoU_{tok}$} & \textbf{$IoU_{agg}$} & \textbf{$F1_{gnd}$} & \textbf{$IoU_{tok}$} & \textbf{$IoU_{agg}$} & \textbf{$F1_{gnd}$} & \textbf{$IoU_{tok}$} & \textbf{$IoU_{agg}$} & \textbf{$F1_{gnd}$} \\
\midrule
% --- InternVL Section ---
\multirow{4}{*}{InternVL2.5-1B} 
 & SFT & \xmark & ---  & 9.56 & 10.73 & 18.67 & 7.81 & 8.60 & 14.91 & 7.93 & 8.83 & 15.14 \\
 & SFT & \cmark & ---  & 9.54 & 10.91 & 19.04 & 8.62 & 9.67 & 16.54 & 8.73 & 9.71 & 16.72 \\
 & RFT & \xmark & GRPO & 9.39 & 10.85 & 18.85 & 8.06 & 8.91 & 15.21 & 8.46 & 9.34 & 15.91 \\
 & RFT & \cmark & GRPO & 10.39 & 11.46 & 19.74 & 9.22 & 9.79 & 16.66 & 8.70 & 9.39 & 16.00 \\
\midrule
% --- Qwen Section ---
\multirow{6}{*}{Qwen2.5-VL-3B} 
 & SFT & \xmark & ---  & 16.42 & 20.04 & 32.50 & 17.38 & 19.75 & 32.05 & 16.12 & 18.33 & 29.96 \\
 & SFT & \cmark & ---  & 15.97 & 19.73 & 31.92 & 17.46 & 19.86 & 32.14 & 16.64 & 18.84 & 30.63 \\
 & RFT & \xmark & GRPO & 16.21 & 18.84 & 30.81 & 16.81 & 18.66 & 30.48 & 15.27 & 17.15 & 28.29 \\
 & RFT & \cmark & GRPO & 16.57 & 19.63 & 31.79 & 17.81 & 19.74 & 31.89 & 17.30 & 19.24 & 31.10 \\
 & RFT & \xmark & GDPO & 16.68 & 19.23 & 31.39 & 17.43 & 19.38 & 31.51 & 15.75 & 17.56 & 28.84 \\
 & RFT & \cmark & GDPO & 16.78 & 19.72 & 31.89 & 18.18 & 20.26 & 32.61 & 18.04 & 20.16 & 32.41 \\
\bottomrule
\end{tabular}%
}
\caption{Quantitative evaluation of visual grounding across different prompt types.}
\label{tab:iou}
\end{table*}

% TODO : decision change rate 가 causal grounding 을 완전히 분석하기에는 불가능하지만, 적어도 causal grounded 되어있다면 유의미한 차이를 보여야함을 명시. 
% \paragraph{Insensitivity to Visual Intervention.} 
\paragraph{Causal Intervention Analysis.}
Previous results suggest that the action distribution shift induced by CF reasoning may be independent of visual evidence. To test this, we conduct a causal intervention study on pre-collision frames, where we remove different visual elements and measure changes in model decisions. Specifically, we consider three masking settings: causal masking (removing the accident-causing object), non-causal masking (removing irrelevant objects), and random masking (removing background regions), and compute the decision change rate by comparing predictions with the original video. As shown in Table~\ref{tab:causal_masking}, human evaluation exhibits the highest sensitivity under causal masking, while non-causal and random masking yield much lower changes, indicating strong reliance on true causal factors. In contrast, most VLMs show no such selectivity, with similar change rates across all masking conditions and this behavior persists regardless of CF reasoning or reinforcement learning. Overall, these results suggest that the action distribution shifts are unlikely to reflect genuine causal reasoning, but are more consistent with visual-agnostic linguistic biases.
% The qualitative results and failure case can be seen in Appendix \ref{appendix:qualitative_masking}.

\begin{table*}[h]
\centering
\footnotesize
\resizebox{\textwidth}{!}{%
\begin{tabular}{@{}l cc cc cccc cccccc cc@{}}
\toprule
\multirow{3}{*}{}
  & \multicolumn{2}{c}{\textbf{InternVL2-1B}}
  & \multicolumn{2}{c}{\textbf{InternVL2-2B}}
  & \multicolumn{4}{c}{\textbf{InternVL2.5-1B}}
  & \multicolumn{6}{c}{\textbf{Qwen2.5-VL-3B}}
  & \multicolumn{2}{c}{\textbf{VideoLLaMA2-7B}} \\
\cmidrule(lr){2-3}
\cmidrule(lr){4-5}
\cmidrule(lr){6-9}
\cmidrule(lr){10-15}
\cmidrule(lr){16-17}
  & \multicolumn{2}{c}{SFT}
  & \multicolumn{2}{c}{SFT}
  & \multicolumn{2}{c}{SFT}
  & \multicolumn{2}{c}{RFT (GRPO)}
  & \multicolumn{2}{c}{SFT}
  & \multicolumn{2}{c}{RFT (GRPO)}
  & \multicolumn{2}{c}{RFT (GDPO)}
  & \multicolumn{2}{c}{SFT} \\
\cmidrule(lr){2-3}
\cmidrule(lr){4-5}
\cmidrule(lr){6-7}
\cmidrule(lr){8-9}
\cmidrule(lr){10-11}
\cmidrule(lr){12-13}
\cmidrule(lr){14-15}
\cmidrule(lr){16-17}
  & \textbf{\xmark} & \textbf{\cmark}
  & \textbf{\xmark} & \textbf{\cmark}
  & \textbf{\xmark} & \textbf{\cmark}
  & \textbf{\xmark} & \textbf{\cmark}
  & \textbf{\xmark} & \textbf{\cmark}
  & \textbf{\xmark} & \textbf{\cmark}
  & \textbf{\xmark} & \textbf{\cmark}
  & \textbf{\xmark} & \textbf{\cmark} \\
\midrule
Causal
  & 0.455 & 0.494
  & 0.408 & 0.451
  & 0.425 & 0.382
  & 0.307 & 0.340
  & 0.457 & 0.524
  & 0.404 & 0.520
  & 0.417 & 0.484
  & 0.434 & 0.508 \\
Non-causal
  & 0.463 & 0.498
  & 0.367 & 0.404
  & 0.416 & 0.373
  & 0.310 & 0.346
  & 0.457 & 0.546
  & 0.416 & 0.504
  & 0.399 & 0.548
  & 0.396 & 0.519 \\
Random
  & 0.475 & 0.508
  & 0.385 & 0.416
  & 0.452 & 0.363
  & 0.295 & 0.355
  & 0.458 & 0.519
  & 0.411 & 0.524
  & 0.388 & 0.531
  & 0.414 & 0.513 \\
\bottomrule
\end{tabular}%
}
\caption{Comparison of causal masking effects across different models and training settings.}
\label{tab:causal_masking}
\end{table*}

\section{Extended Analysis}
\label{sec:more_analysis}
% In summary, SFT 에서 response 생성 시 CF reasoning을 포함하도록 fine-tuning 된 모델들은 action distribution shift를 형성하였다. 이러한 경향은 RFT 이후에도 일관되게 유지되거나 강화되는 것으로 나타났다. 그러나 두 가지 visual grounding 분석에서는 CF reasoning 이 사고 원인 영역에 대한 activation 집중도를 유의미하게 변화시키지 않았으며, 시각적 intervention에 대해서도 decision change rate의 차이가 제한적으로 나타났다. 이는 CF reasoning이 시각적 grounding을 강화하기보다는, 시각적 evidence와 독립적으로 언어적 policy 에 영향을 줄 가능성을 시사한다. 그렇다면 이는 왜 이러한 언어적 편향을 유발하는가? 이 메커니즘을 규명하기 위해, 우리는 SFT 단계에서의 perplexity 변화와 RFT 단계에서의 policy entropy 변화를 분석한다.
\paragraph{Generalization to Trajectory-Level Behavior.} To examine whether the conservative shift persists beyond crash scenarios and under fine-grained trajectory-level evaluation, we conducted an additional study using the existing counterfactual dataset \cite{causal:e2e-vlm:omnidrive}. We compared models trained with and without CF reasoning. For each sample $i$, we compute the maximum lateral displacement of the predicted trajectory, $d_i=\max_t |\ell_{i,t}|$, where $\ell_{i,t}$ is the ego-frame lateral coordinate at time $t$. A trajectory is counted as a lateral maneuver if $d_i \ge \tau$. The lateral maneuver rate over an evaluation subset $\mathcal{S}$ is
\begin{equation}
\rho_\tau(\mathcal{S}) = \frac{1}{|\mathcal{S}|}\sum_{i\in\mathcal{S}} \mathbb{1}\!\left[d_i \ge \tau\right],
\end{equation}
and we report the gap $\Delta_\rho=\rho_\tau^{\text{CF}}-\rho_\tau^{\text{No-CF}}$ for $\tau\in\{0.3, 0.5, 0.8\}$\,m. We evaluate on all samples ($\mathcal{S}_{\text{all}}$) and on the subsets whose ground-truth trajectories contain a lateral maneuver, pooled ($\mathcal{S}_{\text{GT}}$) and split into left-turn ($\mathcal{S}_{\text{L}}$) and right-turn ($\mathcal{S}_{\text{R}}$). As shown in Table~\ref{tab:discussion:omnidrive}, the CF model consistently exhibits lower maneuver rates ($\Delta_\rho<0$) at every threshold, both on $\mathcal{S}_{\text{all}}$ and on $\mathcal{S}_{\text{GT}}$. The same trend holds for both $\mathcal{S}_{\text{L}}$ and $\mathcal{S}_{\text{R}}$, indicating that the conservative shift is not merely attributable to the prevalence of straight-driving scenarios. These results show that the conservative behavioral bias induced by CF reasoning is also observed in general driving scenarios and under low-level trajectory evaluation. Moreover, performance improvements on nuScenes do not necessarily reflect improved reasoning: since straight trajectories make up most of the ground truth in the dataset~\cite{causal:e2e:ad-mlp,ad:unitraj}, conservative predictions can achieve low average error, which is consistent with the model trained with CF reasoning.



\begin{table*}[h]
\centering
\setlength{\tabcolsep}{5pt}
\begin{tabular}{@{}clrrrr@{}}
\toprule
$\tau$ (m) & Subset & $N$ & $\rho_\tau^{\text{No-CF}}$ & $\rho_\tau^{\text{CF}}$ & $\Delta_\rho$  \\
\midrule
\multirow{4}{*}{0.3m}
& $\mathcal{S}_{\text{all}}$& 5,119 & 8.85\%  & 3.07\% & $-5.78$ pp  \\
& $\mathcal{S}_{\text{GT}}$ & 2,286 & 17.94\% & 5.95\% & $-11.99$ pp \\
& $\mathcal{S}_{\text{L}}$  & 874   & 29.52\% & 11.21\% & $-18.31$ pp \\
& $\mathcal{S}_{\text{R}}$  & 1,412 & 10.76\% & 2.69\% & $-8.07$ pp \\
\midrule
\multirow{4}{*}{0.5m}
& $\mathcal{S}_{\text{all}}$& 5,119 & 6.58\%  & 1.93\% & $-4.65$ pp \\
& $\mathcal{S}_{\text{GT}}$ & 1,676 & 18.85\% & 5.55\% & $-13.31$ \\
& $\mathcal{S}_{\text{L}}$  & 697   & 32.14\% & 11.19\% & $-20.95$ pp \\
& $\mathcal{S}_{\text{R}}$  & 979   & 9.40\%  & 1.53\% & $-7.87$ pp \\
\midrule
\multirow{4}{*}{0.8m}
& $\mathcal{S}_{\text{all}}$& 5,119 & 3.48\%  & 0.84\% & $-2.64$ pp \\
& $\mathcal{S}_{\text{GT}}$ & 1,244 & 13.34\% & 3.05\% & $-10.29$ pp \\
& $\mathcal{S}_{\text{L}}$  & 512   & 25.39\% & 6.64\% & $-18.75$ pp \\
& $\mathcal{S}_{\text{R}}$  & 732   & 4.92\%  & 0.55\% & $-4.37$ pp \\
\bottomrule
\end{tabular}
\caption{Lateral maneuver rate $\rho_\tau$ of No-CF and CF models on OmniDrive, where a trajectory counts as a maneuver if its maximum lateral displacement is at least $\tau$. $\Delta_\rho=\rho_\tau^{\text{CF}}-\rho_\tau^{\text{No-CF}}$. $\mathcal{S}_{\text{all}}$: all samples; $\mathcal{S}_{\text{GT}}$: samples whose ground-truth trajectory contains a lateral maneuver; $\mathcal{S}_{\text{L}}$, $\mathcal{S}_{\text{R}}$: its left- and right-turn subsets ($\mathcal{S}_{\text{GT}}=\mathcal{S}_{\text{L}}\cup\mathcal{S}_{\text{R}}$).}
\label{tab:discussion:omnidrive}
\end{table*}

\paragraph{Investigating Semantic Priors as a Cause of the Action Shift.} Furthermore, to investigate whether the conservative shift arises from hazard-specific linguistic priors rather than visual and causal grounding, we conducted a conflict test. We injected either adversarial hazard-related text or safety-neutral text into normal driving videos without accidents and measured the false positive rate (FPR) of \texttt{Stop} and \texttt{Decel.} predictions. We define *hazard specificity* as the difference in FPR between hazard-related and safety-neutral text injections, where a positive value indicates a stronger response to hazard-related text. As shown in Table~\ref{tab:discussion:semantic}, models with CF reasoning already predict \texttt{Stop} and \texttt{Decel.} in 61--69\% of cases, compared with 12--25\% for models without CF reasoning, confirming their conservative bias. Under text injection, No-CF models exhibit near-zero or weakly positive hazard specificity ($+0.71$ to $+5.00$\,pp), whereas CF models exhibit negative hazard specificity across all three training settings ($-6.50$ to $-14.21$\,pp), indicating that they respond more strongly to safety-neutral text than to hazard-related text. This pattern runs counter to what would be expected if hazard-specific semantic priors were driving the conservative shift. Nevertheless, the elevated FPR of CF models under both types of injection suggests that their predictions are sensitive to injected text regardless of its safety-related content. Taken together, these results suggest that hazard-related semantic priors alone cannot account for the observed conservative shift.


\begin{table*}[h]
\centering
\footnotesize
\resizebox{\textwidth}{!}{%
\begin{tabular}{@{}l cc cc cc@{}}
\toprule
\multirow{2}{*}{\textbf{Metric}}
  & \multicolumn{2}{c}{\textbf{Qwen2.5-VL-3B (SFT)}}
  & \multicolumn{2}{c}{\textbf{Qwen2.5-VL-3B (GRPO)}}
  & \multicolumn{2}{c}{\textbf{Qwen2.5-VL-3B (GDPO)}} \\
\cmidrule(lr){2-3}
\cmidrule(lr){4-5}
\cmidrule(lr){6-7}
  & \textbf{No-CF} & \textbf{CF}
  & \textbf{No-CF} & \textbf{CF}
  & \textbf{No-CF} & \textbf{CF} \\
\midrule
Baseline FPR
  & 25.00\% & 69.00\%
  & 21.58\% & 66.96\%
  & 12.35\% & 61.01\% \\

Adversarial FPR
  & 33.25\% & 72.25\%
  & 31.21\% & 63.24\%
  & 22.92\% & 60.20\% \\

Neutral FPR
  & 28.25\% & 78.75\%
  & 29.65\% & 77.46\%
  & 22.21\% & 70.32\% \\

Hazard Specificity (Adv. - Neu.)
  & +5.00 pp & $-$6.50 pp
  & +1.56 pp & $-$14.21 pp
  & +0.71 pp & $-$10.12 pp \\
\bottomrule
\end{tabular}%
}
\caption{False positive rates and hazard specificity under baseline,
adversarial, and neutral conditions across different model architectures
and training methods. No-CF and CF denote training without and with
CF reasoning, respectively. Hazard specificity is defined as
$\mathrm{FPR}_{\mathrm{Adv}}-\mathrm{FPR}_{\mathrm{Neu}}$.}
\label{tab:discussion:semantic}
\end{table*}




% TODO : 기존 연구들이 무얼 놓쳤는가
% TODO : 본 연구가 왜 필요한가
\section{Related Work}
% Recent studies have begun to critically examine the "black-box" nature of end-to-end autonomous driving models, revealing that high performance often stems from shortcut learning rather than genuine scene understanding. \cite{dataset:drivebench} 는 blah blah...  \cite{e2e-ad:is-ego-status} first exposed this issue by demonstrating that simple MLP-based policies relying solely on ego-status (e.g., speed, acceleration) could match the performance of complex vision-based models, suggesting a heavy reliance on inertia priors rather than visual perception. Expanding this scrutiny to VLMs, \cite{causal:e2e-vlm:more_than_meets_the_eye} uncovered a disturbing disconnect between reasoning and planning; they showed that while VLMs generate plausible Chain-of-Thought (CoT) explanations, their actual planning decisions remain dependent on ego-status shortcuts, rendering the reasoning process superficial, \textbf{but their study was limited to 주행 성능으로만 reasoning capability 를 평가하려했다.} In contrast, vision part 에 대한 명시적 분석과 reasoning capacity 를 기반으로 해당 현상을 분석했다.
% Recent studies have highlighted various limitations faced by VLMs in autonomous driving. For example, \cite{dataset:drivebench} emphasized the reliability of visual grounding in complex real-world environments, showing that VLMs tend to produce plausible but unsupported responses under severe visual degradation, raising concerns about their trustworthiness in safety-critical decision-making. Additionally, the previously noted issue of excessive reliance on ego and navigation status in end-to-end autonomous driving \cite{e2e-ad:is-ego-status} has been reported to persist in VLM-based approaches as well \cite{causal:e2e-vlm:more_than_meets_the_eye}. Independently of these limitations, recent work has introduced CF reasoning \cite{causal:e2e-vlm:omnidrive, causal:e2e-vlm:simlingo, causal:cf_vla}, demonstrating meaningful improvements in safety-relevant driving performance. However, these studies offer limited analysis of how CF reasoning induces structural changes in the model’s internal representations or decision-making processes. In particular, the relationship between CF reasoning and previously identified issues such as modality misalignment and linguistic priors remains underexplored. Motivated by this gap, this study seeks to reinterpret the effects of CF reasoning from a more fundamental perspective by analyzing changes in the model’s uncertainty, exploration tendencies, and action biases, while highlighting both its structural limitations and potential benefits. More comprehensive review of related work is provided in Section \ref{appendix:discussion}.

\paragraph{Shortcut Learning and Causal Confusion.}
% 분포 변화(distribution shift) 상황에서도 안정적인 성능을 보장하는 것은 자율주행에서 핵심적인 문제이며, 그중에서도 causal reasoning은 spurious correlation보다 환경 변화에 대해 더 invariant하게 유지된다는 점에서 특히 강건한 접근으로 간주된다. 반면 imitation learning 기반 접근법은 전문가의 행동을 모방하는 방식으로 학습되기 때문에 shortcut learning에 취약하다 \cite{causal:e2e:causal_confusion, causal:e2e:driveadapter, causal:e2e:exploring_causality_in_e2e}. 이는 모델이 실제 인과 구조가 아닌 데이터 내 통계적 상관관계에 의존하여 의사결정을 학습하는 현상으로, causal confusion이라고도 불린다. 이러한 현상은 모델이 본질적인 인과 메커니즘이 아니라 표면적 패턴에 의존하게 만드는 문제를 초래한다. 대표적인 사례로, 시각 정보 없이 ego-status와 navigation information만으로도 경쟁력 있는 주행 성능이 가능하다는 결과가 보고된 바 있다 \cite{causal:e2e:ad-mlp, e2e-ad:is-ego-status}. 이러한 문제를 완화하기 위해 causal factor를 명시적으로 분리하려는 연구들이 제안되었으나 \cite{causal:e2e:cadet, causal:e2e:causalvad, causal:e2e:decoupling_perception_and_ego_status}, 복잡한 장면 이해 및 문맥적 추론에서는 여전히 한계를 보이고 있다.
Ensuring robust performance under distribution shift is a central challenge in autonomous driving. Thus, causal reasoning is regarded as a particularly promising solution due to its invariance to environmental changes compared to spurious correlations. However, existing works based on imitation learning, which learn by mimicking expert demonstrations, are inherently vulnerable to shortcut learning \cite{causal:e2e:causal_confusion, causal:e2e:driveadapter, causal:e2e:exploring_causality_in_e2e}. This phenomenon, also referred to as causal confusion, arises when models rely on statistical correlations in the data rather than underlying causal structure. As a result, models tend to depend on superficial patterns instead of causal reasoning. A representative example is the observation that competitive driving performance can be achieved even using only ego-status and navigation information \cite{causal:e2e:ad-mlp, e2e-ad:is-ego-status}. To mitigate such issues, several studies have proposed explicitly disentangling causal factors \cite{causal:e2e:cadet, causal:e2e:causalvad, causal:e2e:decoupling_perception_and_ego_status}. However, these approaches still struggle in complex scene understanding and contextual reasoning scenarios, indicating that shortcut learning remains an open challenge in end-to-end autonomous driving.


\paragraph{Illusion of Causal Reasoning in VLM for Autonomous Driving.}
% 이러한 한계를 극복하기 위해 최근에는 VLM을 활용한 자율주행 연구가 등장하였다. 이는 복잡한 주행 시나리오 context를 이해하는데 있어서 language modality를 통해 coarse semantic representation을 활용하며, 기존 end-to-end 방식 대비 closed-loop 환경에서 더 높은 성능을 보이는 것으로 보고되었다 \cite{causal:e2e-vlm:simlingo}. 이후 연구들은 chain-of-thought \cite{e2e-vlm:drivelm}, chain-of-causation \cite{causal:e2e-vlm:coc-vla}, counterfactual reasoning \cite{causal:e2e-vlm:omnidrive, causal:cf_vla} 등을 decision-making 이전 단계의 intermediate reasoning으로 도입하여 성능을 추가적으로 향상시키고자 하였다. 이에 따른 성능 향상은 언어 기반 reasoning이 마치 causal reasoning을 실제로 수행하는 것처럼 보이게 만들며, 그 결과 VLM이 언어적 수준에서 인과적 이해를 획득했다는 암묵적인 해석으로 이어지고 있다.

To address limitations of conventional end-to-end approaches, recent studies have explored VLMs for autonomous driving, leveraging language as a coarse semantic representation to better capture complex scene contexts and achieving improved performance in closed-loop settings \cite{causal:e2e-vlm:simlingo, e2e-vlm:orion, e2e-vlm:autovla}. Building on this, prior work has introduced intermediate linguistic reasoning steps such as chain-of-thought \cite{e2e-vlm:drivelm}, chain-of-causation \cite{causal:e2e-vlm:coc-vla}, and counterfactual reasoning \cite{causal:cf_vla}, leading to consistent performance gains and an implicit assumption that such reasoning reflects genuine causal understanding. However, it remains unclear whether these improvements stem from causal reasoning or from shortcut learning induced by the language modality, as standard driving metrics are confounded and insufficient for analyzing internal reasoning processes \cite{causal:e2e-vlm:more_than_meets_the_eye}. Recent evidence from CausalDriveBench \cite{causal:e2e-vlm:causaldrivebench} further highlights this limitation, showing that causal QA accuracy is uncorrelated with trajectory prediction accuracy across driving VLAs, indicating that conventional driving metrics alone cannot assess whether a model's linguistic reasoning reflects genuine causal understanding. To address this, we introduce a framework to evaluate the behavioral effects of CF reasoning by quantifying the induced shifts in a discrete meta-action space and validating whether these shifts persist under reinforcement learning, thereby ruling out training artifacts. This enables a systematic analysis of the gap between linguistic shortcut and genuine causal understanding.



\section{Conclusion and Limitations}
\label{sec:conclusion}
% TODO (CONCLUDE) : 배경, 요약, 등등
In this study, we investigate the effect of counterfactual (CF) reasoning in VLMs for autonomous driving. Our results indicate that CF reasoning consistently induces an action distribution shift, and that these changes primarily originate from language-level reasoning rather than genuine causal reasoning grounded in visual observations. This finding suggests that CF reasoning can guide the model in determining what to say, but does not necessarily teach the model how to reason. In other words, the behavior shift is more plausibly attributed to the learning of a language-induced spurious correlation, rather than the internalization of genuine causal reasoning.

However, this study focuses on open-source VLMs that allow access to internal probability distributions and learning dynamics. Consequently, the generalizability of our findings to closed-source or specialized VLMs \cite{causal:e2e-vlm:alpamayo, causal:cf_vla} remains uncertain, as our analyses based on perplexity and policy entropy require access to internal model probabilities. Therefore, our conclusions are limited to settings where such internal signals are observable, highlighting the need for future work to extend CF reasoning analysis beyond these constraints. We further discuss potential causes of the observed action distribution shift and outline future directions in Section \ref{appendix:discussion}.

\section{Acknowledgement}
This work was supported by the National Research Foundation of Korea (NRF) grant funded by the Korea government (MSIT) (RS-2026-25510069), the Institute of Information and Communications Technology Planning and Evaluation (IITP) grant funded by the Korea government (No. 2022-0-00469), and, in part, by the BK21 FOUR (Connected AI Education \& Research Program for Industry and Society Innovation, KAIST EE, No. 4120200113769).

\clearpage
\bibliographystyle{unsrt}
% \bibliographystyle{plainnat} 
\bibliography{ref}



%%%%%%%%%%%%%%%%%%%%%%%%%%%%%%%%%%%%%%%%%%%%%%%%%%%%%%%%%%%%


\clearpage
\appendix

\renewcommand{\thesection}{\Alph{section}} % Optional: Keep if you liked this numbering

% --- titletoc Appendix ToC Start ---
\startcontents[appendix] % Start collecting info for a ToC named 'appendix'
\section*{Appendix} % Add an unnumbered title for this ToC
\printcontents[appendix] % Print the ToC named 'appendix'
               {}{1} % Format: {} = no prefix, {} = default formatting, 1 = start at level 1 (section)
               {\setcounter{tocdepth}{2}} % Set depth to 1 (only sections). Use 2 for subsections etc.

\clearpage
% TODO : 본 연구에서 발견한 결과와 더불어, 선행 연구들을 기반으로 후속 연구 방향 제시.
% 이 논문이 '분석'에 초점을 맞추고 있으므로, 논문에서 보인 한계나 분석 결과와 상응하는 INSIGHT 를 줘야함 

\section{Discussion}
\label{appendix:discussion}
% 왜 중요한가, 무엇을 의미하는가, 어디까지 주장할 수 있는가, 다음에는 무엇을 할 수 있는가
\paragraph{Why Does CF Reasoning Induce Behavioral Shifts Without Visual Grounding?} A key factor is that CF reasoning in current VLM-based autonomous driving systems is primarily learned at the language level. During supervised fine-tuning, supervision is provided only in the form of textual reasoning outputs, without explicit alignment between these texts and the corresponding visual evidence. As a result, CF reasoning is optimized as a language modeling objective rather than a visually grounded causal reasoning, leading to a natural decoupling between linguistic reasoning and visual causal structure. This interpretation is consistent with our observations of invariant attention patterns and non-selective responses under causal intervention. These findings further challenge the implicit assumption that performance improvements from language-based reasoning necessarily reflect genuine causal understanding. In practice, both LLMs and VLMs exhibit fundamental limitations in stable causal inference, particularly under distribution shift in long-tail scenarios \cite{causal:whatif_tv_off?, causal:llm_cannot_discover_causality, dataset:fail2drive}. 

\paragraph{Is Visual Grounding Sufficient for Causal Decision-Making?} To address the limitation, prior work has proposed leveraging counterfactual consistency to encourage reliance on visual evidence rather than purely linguistic patterns \cite{causal:defacto} and further explored integrating LLM-based agents with causal world models, extending decision-making beyond next-token prediction toward structured environments that support intervention and counterfactual reasoning \cite{causal:bridging_llms_and_causal_world_model}. These directions highlight the importance of developing CF reasoning as a learning signal that is aligned with physical grounding, rather than relying solely on linguistic supervision. However, visual grounding alone does not necessarily ensure that the resulting reasoning is effectively reflected in downstream decision-making. In other words, establishing a correspondence between reasoning and visual evidence is distinct from ensuring that such reasoning appropriately guides the model's actions \cite{causal:e2e-vlm:grounding_shapes_plan,robotics:linking_causal_to_gen_in_vla}.

\paragraph{Future Directions.}
% (1) VLM 이 자율주행에 있어서 language modality 를 통해 인과 추론을 가능하게 만든다는 주장과 더불어 향상된 driving score 가 결국 causal understanding 에서 기인했다는 암묵적인 주장을 배제하려면 평가 방식을 바꿔야함. casual understanding 과 driving 이라는 downstream task 와의 gap 을 평가하는 방식이 필요.  
% (2) 인과추론을 학습시키는것이 방대하고 어려운 문제라면, 적어도 shorcut 을 저해하는 방향으로 우회할 수 있다
Despite their ability to capture rich correlational structures and derive coarse semantic representations of scenes in autonomous driving through large-scale pretraining \cite{causal:e2e-vlm:from_representational}, VLMs still lack reliable causal reasoning capabilities. Therefore, the implicit assumption that language-based reasoning enables VLMs to perform causal reasoning in autonomous driving, and that improvements in driving performance necessarily reflect genuine causal understanding, warrants reconsideration. This calls for evaluation paradigms that go beyond conventional driving performance metrics to explicitly examine the gap between causal understanding and downstream driving behavior. In particular, evaluating the alignment between reasoning and action could provide insights into whether the model's reasoning genuinely informs its decisions, rather than merely accompanying behavioral improvements.

Furthermore, given the substantial challenges involved in endowing VLMs with explicit causal reasoning capabilities, an alternative research direction is to pursue more tractable objectives that mitigate the influence of causal confounders, such as linguistic shortcuts. Rather than directly solving the broader problem of causal reasoning, identifying and suppressing such shortcuts may offer a practical step toward more reliable and causally grounded decision-making \cite{causal:e2e:causalvad, causal:e2e:cadet, causal:e2e:decoupling_perception_and_ego_status}.  



% \clearpage
\section{Detailed Implementation}
\label{appendix:implementations}
\subsection{Group Relative Policy Optimization (GRPO)}
\label{appendix:implementations_grpo}
For each query $q$ (comprising an input image and prompt), GRPO samples a group of $G$ outputs $\{o_1, o_2, \dots, o_G\}$ from the old policy $\pi_{\theta_{\text{old}}}$. It then calculates the reward $r_i$ for each output and updates the policy using relative advantages computed based on the group mean. The objective function is defined as:

$$\begin{aligned}
\mathcal{J}_{\mathrm{GRPO}}(\theta) &= \mathbb{E}_{q \sim P(Q), \{o_i\}_{i=1}^G \sim \pi_{\theta_{\text{old}}}} \left[ \frac{1}{G} \sum_{i=1}^G \mathcal{J}_i - \beta \mathbb{D}_{\mathrm{KL}}(\pi_\theta \| \pi_{\text{ref}}) \right], \\
\mathcal{J}_i &= \min \left( w_i \hat{A}_i, \text{clip}(w_i, 1-\epsilon, 1+\epsilon) \hat{A}_i \right),
\end{aligned}$$

where $w_i = \frac{\pi_\theta(o_i|q)}{\pi_{\theta_{\text{old}}}(o_i|q)}$ is the importance sampling ratio, and $\hat{A}_i$ denotes the group-relative advantage.

\subsection{Group reward-Decoupled Normalization Policy Optimization (GDPO)}
\label{appendix:implementations_gdpo}
Unlike GRPO, which directly incorporates group-relative advantages into the policy objective, GDPO decouples group-wise reward normalization from policy optimization. Specifically, the normalized rewards $\tilde{r}_i$ are treated as fixed advantages, and no gradients are propagated through the normalization statistics $\mu_q$ and $\sigma_q$. The GDPO objective is formulated in a PPO-style manner as:

$$\begin{aligned}
\mathcal{J}_{\mathrm{GDPO}}(\theta) &= \mathbb{E}_{q \sim P(Q), \{o_i\}_{i=1}^G \sim \pi_{\theta_{\text{old}}}} \left[ \frac{1}{G} \sum_{i=1}^G \mathcal{L}_i^{\mathrm{GDPO}} - \beta \mathbb{D}_{\mathrm{KL}}(\pi_\theta \| \pi_{\text{ref}}) \right], \\
\mathcal{L}_i^{\mathrm{GDPO}} &= \min \left( w_i \tilde{r}_i, \text{clip}(w_i, 1-\epsilon, 1+\epsilon)\tilde{r}_i \right),
\end{aligned}$$

where $w_i = \frac{\pi_\theta(o_i|q)}{\pi_{\theta_{\text{old}}}(o_i|q)}$ denotes the same importance sampling ratio used in GRPO.


% TODO : "does RLVR... " 논문이랑 pass@k 구한 방법이 다름. estimator 라는 단어 참고 
\subsection{Pass@k}
\label{appendix:implementations_passk}
\begin{equation}
\text{pass}@k := \mathbb{E} \left[ 1 - \frac{\binom{n-c}{k}}{\binom{n}{k}} \right]
\label{eq:pass_at_k}
\end{equation}
\noindent where $n$ is the total number of generated samples, and $c$ denotes the count of correct (safe) outputs. This metric provides a more robust evaluation of the model's potential than a single-shot prediction, as it accounts for the probabilistic nature of the model's reasoning.



\subsection{Perplexity}
\label{appendix:implementations_perplexity}
Given a model $m$, a prompt $x$, and a sequence of tokens $Y = (y_1, \dots, y_T)$, the perplexity is defined as the exponentiated average negative log-likelihood:

$$
\mathrm{PPL}_m(Y \mid x) = \exp \left( -\frac{1}{T} \sum_{t=1}^{T} \log P_m(y_t \mid x, y_1, \ldots, y_{t-1}) \right)
$$



\subsection{Policy Entropy}
\label{appendix:implementations_policy_entropy}
\begin{equation}
    H_t(\pi_\theta) = -\sum_{v \in \mathcal{V}} \pi_\theta(v \mid \mathbf{x}, \mathbf{y}_{<t}) \log \pi_\theta(v \mid \mathbf{x}, \mathbf{y}_{<t})
\end{equation}

where $\mathcal{V}$ is the vocabulary. To ensure a fair comparison across models with varying response lengths (e.g., concise standard SFT vs. verbose Counterfactual models), we employ a \textit{length-normalized} average entropy:

\begin{equation}
    \mathcal{H}_{\text{policy}}(\pi_\theta, \mathbf{x}, \mathbf{y}) = \frac{1}{T} \sum_{t=1}^{T} H_t(\pi_\theta)
    \label{eq:policy_entropy}
\end{equation}

% TODO : policy entropy 측정 방식의 공정성 및 response length 로 인한 confound  
By normalizing for sequence length $T$, $\mathcal{H}_{\text{policy}}$ reduces the influence of verbosity on entropy comparisons. We note, however, that this normalization may not fully eliminate length-related effects, as per-token entropy tends to decrease toward the end of longer sequences (e.g., during repetitive closing patterns). Consequently, if CF models produce longer responses on average, the length-normalized estimator is likely to \textit{underestimate} their true entropy relative to shorter non-CF responses. This structural bias works \textit{against} our hypothesis, making the observed entropy elevation in CF models a conservative estimate of the underlying distributional shift. 

% TODO : 
To further validate that the entropy elevation is not primarily driven by response length, we compute the Pearson correlation coefficient between mean completion length and policy entropy during RFT. As shown in Table~\ref{tab:entropy_corr}, the correlation is consistently lower for CF models (0.365--0.392) than for non-CF models (0.439--0.499) across all settings, suggesting that response length explains less of the entropy variance in CF models than in their non-CF counterparts.
\begin{table*}[h]
\centering
\begin{tabular}{lcc}
\toprule
\textbf{Model} & \textbf{w/ CF} & \textbf{w/o CF} \\
\midrule
InternVL2.5-1B (GRPO) & 0.365 & 0.499 \\
Qwen2.5-VL-3B (GRPO)  & 0.392 & 0.461 \\
Qwen2.5-VL-3B (GDPO)  & 0.370 & 0.439 \\
\bottomrule
\end{tabular}
\caption{Pearson correlation coefficient between mean completion length and policy entropy during RFT. CF models consistently exhibit lower correlation than their non-CF counterparts across all settings, suggesting that the entropy elevation in CF models is not primarily driven by response length.}
\label{tab:entropy_corr}
\end{table*}

\subsection{Metrics for Token Activation Map (TAM) Evaluation}
\label{appendix:implementations_iou_metrics}

To quantitatively evaluate whether counterfactual reasoning 
enhances visual grounding on accident-causing objects, we 
employ Token Activation Map (TAM)~\cite{attention:tam}, 
which attributes each output token's generation to specific 
visual regions. Unlike standard attention visualization that 
conflates multi-head and multi-layer interactions, TAM 
isolates the contribution of vision tokens to each decoded 
token through targeted perturbation of hidden states.


For a generated token $t_i$ at decoding step $i$, TAM 
computes the activation map $\mathbf{A}_i \in \mathbb{R}^{H \times W}$ 
by measuring the sensitivity of the token's logit to each 
vision token embedding. Specifically, given vision tokens 
$\{\mathbf{v}_j\}_{j=1}^{H \times W}$ from the vision encoder 
and the LLM hidden state $\mathbf{h}_i$ before the language 
head, TAM perturbs each vision token and observes the change 
in output probability:

\begin{equation}
\mathbf{A}_i(x, y) = \left| \log p(t_i | \mathbf{h}_i) - 
\log p(t_i | \mathbf{h}_i^{(x,y)\to 0}) \right|
\end{equation}

where $\mathbf{h}_i^{(x,y)\to 0}$ denotes the hidden state 
when vision token at spatial position $(x, y)$ is zeroed out. 
The resulting map $\mathbf{A}_i$ is then normalized to 
$[0, 1]$ and resized to the original image resolution for 
bounding box comparison. Based on it, we define three complementary metrics to measure visual grounding at different granularities:

\paragraph{Token-level IoU ($IoU_{tok}$)}
For each causal keyword token (e.g., ``car'', ``pedestrian'') 
in the generated response, we compute the Intersection over 
Union (IoU) between its activation map $\mathbf{A}_i$ and 
the ground-truth bounding box $\mathcal{B}_{causal}$ of the 
accident-causing object:

\begin{equation}
IoU_i = \frac{|\mathbf{M}_i \cap \mathcal{B}_{causal}|}
{|\mathbf{M}_i \cup \mathcal{B}_{causal}|}
\end{equation}

where $\mathbf{M}_i = \{\mathbf{A}_i \geq 0.2 \cdot \max(\mathbf{A}_i)\}$ 
is the binary mask obtained by thresholding at 20\% of the 
peak activation, following standard practice in CAM-based 
methods~\cite{attention:grad-cam}. Token-level IoU is then 
averaged over all causal tokens:

\begin{equation}
IoU_{tok} = \frac{1}{N_{causal}} \sum_{i \in \mathcal{T}_{causal}} IoU_i
\end{equation}

where $\mathcal{T}_{causal}$ is the set of causal keyword 
token indices and $N_{causal} = |\mathcal{T}_{causal}|$.

\paragraph{Aggregated IoU ($IoU_{agg}$)}
To capture the cumulative visual focus across the entire 
reasoning process, we aggregate activation maps from all 
causal tokens by summing their normalized maps:

\begin{equation}
\mathbf{A}_{agg} = \sum_{i \in \mathcal{T}_{causal}} 
\frac{\mathbf{A}_i - \min(\mathbf{A}_i)}
{\max(\mathbf{A}_i) - \min(\mathbf{A}_i)}
\end{equation}

The aggregated IoU is then computed on the combined map:

\begin{equation}
IoU_{agg} = \frac{|\mathbf{M}_{agg} \cap \mathcal{B}_{causal}|}
{|\mathbf{M}_{agg} \cup \mathcal{B}_{causal}|}
\end{equation}

where $\mathbf{M}_{agg} = \{\mathbf{A}_{agg} \geq 0.2 \cdot 
\max(\mathbf{A}_{agg})\}$. This metric reflects whether the 
model's attention, when viewed holistically across all 
reasoning tokens, aligns with the causal object.

\paragraph{Grounding F1 ($F1_{gnd}$)}
IoU alone may underweight scenarios where the model attends 
to the causal region but also includes spurious areas, or 
conversely, misses parts of the object. We therefore compute 
precision and recall on the aggregated binary mask:

\begin{align}
\text{Precision} &= \frac{|\mathbf{M}_{agg} \cap \mathcal{B}_{causal}|}
{|\mathbf{M}_{agg}|} \\
\text{Recall} &= \frac{|\mathbf{M}_{agg} \cap \mathcal{B}_{causal}|}
{|\mathcal{B}_{causal}|}
\end{align}

The Grounding F1 score balances these two aspects:

\begin{equation}
F1_{gnd} = \frac{2 \cdot \text{Precision} \cdot \text{Recall}}
{\text{Precision} + \text{Recall}}
\end{equation}

A high $F1_{gnd}$ indicates that the model's visual attention 
is both concentrated on the causal object (high precision) 
and covers it comprehensively (high recall), suggesting 
genuine visual grounding rather than diffuse or misaligned 
activation.

In summary, if counterfactual reasoning improves causal understanding, CF-enabled models are expected to exhibit higher $IoU_{tok}$, indicating that each reasoning token attends more precisely to the accident cause; higher $IoU_{agg}$, reflecting that the cumulative reasoning process focuses more consistently on the causal region; and higher $F1_{gnd}$, suggesting that the attention is both precise (i.e., low false positives) and comprehensive (i.e., low false negatives).

However, as shown in Table~\ref{tab:iou}, no such improvements 
are observed after instruction fine-tuning with counterfactual reasoning, indicating that the decision shifts via counterfactual reasoning operate independently of visual evidence.


% \clearpage
% \section{More Related Works and Discussions}
% \label{appendix:related_worsk_and_discussions}

% \paragraph{Driving with Language.}
% % 최근 자율주행 분야에서는 언어를 활용하여 주행 의사결정을 보조하거나 직접 생성하려는 연구들이 활발히 이루어지고 있다. 초기 연구들은 perception module의 출력을 자연어로 설명하거나, 주행 상황에 대한 질의응답을 수행하는 방식에 초점을 맞추었으나, Vision-Language Models (VLMs)의 발전과 함께 언어를 의사결정 과정의 핵심 요소로 활용하려는 방향으로 확장되고 있다. 특히, 최근의 end-to-end 자율주행 프레임워크에서는 언어 instruction이나 reasoning을 기반으로 주행 정책을 직접 생성하는 접근이 제안되고 있으며, 이는 모델이 단순한 행동 모방을 넘어 상황에 대한 해석과 의도를 반영한 의사결정을 수행할 수 있도록 한다는 점에서 주목받고 있다. Chain-of-Thought (CoT)에서 나아가, Chain-of-Causation (CoC) 와 counterfactual reasoning dataset 은  causal reasoning을 유도하려는 시도로 이루어지고 있다.
% Recent advances in autonomous driving have increasingly explored the use of language to assist or directly generate driving decisions. Early approaches primarily focused on describing the outputs of perception modules in natural language or performing question answering about driving scenarios. Furthermore, with the development of VLMs, there has been a shift toward leveraging language as a core component of the decision-making process. In particular, recent end-to-end autonomous driving frameworks propose generating driving policies directly based on language instructions or reasoning, enabling models to move beyond simple behavior imitation and instead make decisions that reflect contextual understanding and intent. Building upon Chain-of-Thought (CoT), approaches such as Chain-of-Causation (CoC) and counterfactual (CF) reasoning datasets have been introduced to encourage causal reasoning.


% \paragraph{Limitations of VLMs and Post-Training.} 
% However, these approaches still suffer from a fundamental limitation, namely modality misalignment arising from reliance on linguistic priors. This can lead to plausible yet unsupported reasoning that is not grounded in actual observations, particularly when visual information is incomplete or ambiguous. Such issues may manifest as decision biases in VLM-based autonomous driving and robotics systems. To mitigate these limitations, post-training techniques such as reinforcement learning with human feedback (RLHF) and reinforcement learning with verifiable rewards (RLVR) have been introduced. However, these methods primarily operate by reweighting the policy within the distribution shaped by pretrained priors, and thus struggle to effectively explore novel reasoning patterns beyond this prior. In particular, in the vast action space of language generation, samples that deviate from the prior are likely to be implausible or receive low rewards, causing the training process to converge toward reinforcing the existing prior. As a result, these post-training methods tend to restructure rather than expand the model’s reasoning capacity. This raises a fundamental question of whether language-based reasoning reflects causal understanding of the environment, or merely amounts to a redistribution of linguistic priors.
% % 그러나 이러한 접근은 언어 priors에 대한 의존으로 인한 modality misalignment 라는 근본적인 한계를 여전히 내포하고 있다. 이는 시각적 정보가 불완전하거나 모호한 상황에서 실제 관측과 무관한 그럴듯한 추론을 생성하는 문제로 이어질 수 있다. 이는 VLM 기반 자율주행 및 로보틱스 연구에서 decision bias로 나타날 수 있다. 이러한 한계를 완화하기 위해 reinforcement learning with human feedback (RLHF)나 reinforcement learning with verifiable rewards (RLVR)와 같은 post-training 기법들이 도입되고 있다. 그러나 이러한 방법들은 pretrained prior에 의해 형성된 행동 분포 내에서 정책을 재조정하는 데 그치며, prior를 벗어나는 새로운 추론 패턴을 효과적으로 탐색하지 못한다는 한계를 가진다. 특히, 언어 생성의 vast action space에서는 prior 밖의 샘플이 비합리적이거나 낮은 보상을 받을 가능성이 높기 때문에, 학습 과정이 기존 prior를 더욱 강화하는 방향으로 수렴하게 된다. 그 결과, 이러한 post-training 기법들은 모델의 추론 능력을 근본적으로 확장하기보다는 기존 능력을 재구성하는 수준에 머무르게 된다. 이에 따라, 언어 기반 reasoning이 실제 환경에 대한 인과적 이해를 반영하는지, 혹은 단순히 언어적 분포의 재조정에 그치는지에 대한 근본적인 의문이 제기된다.



% % future work가 “해결책”처럼 보이면 실패고, “새로운 문제 정의 (패러다임)”처럼 보이면 성공. 즉, open-question 같은 "열린 결말"로 남겨놓아야 한다. 
% \paragraph{Discussion : Possible Future Works.}
% % 본 연구는 이러한 문제의식을 바탕으로, CF reasoning이 모델의 행동에 미치는 영향을 분석하였다. 그 결과, CF reasoning은 불확실성과 탐색 특성을 변화시켜 행동 분포를 재구성하지만, 이러한 변화는 시각적 근거에 기반한 causal grounding이라기보다 언어적 prior의 재편성에 의해 설명된다. 이러한 관점에서, 향후 연구는 CF reasoning이 실제 물리적 환경과 얼마나 일관되게 연결될 수 있는지, 즉 modality alignment를 보다 근본적으로 재고할 필요가 있다. 나아가, 단순한 counterfactual supervision을 넘어 모델의 reasoning capacity 자체를 확장할 수 있는 학습 패러다임이 요구된다. 이는 언어적 추론과 물리적 세계 모델 간의 정합성을 어떻게 확보할 것인지, 그리고 prior에 의해 제한된 탐색을 어떻게 넘어설 것인지에 대한 새로운 연구 질문으로 이어진다.
% Building on this motivation, this study analyzes the impact of CF reasoning on model behavior. Our findings show that CF reasoning reshapes the action distribution by altering uncertainty and exploration characteristics; however, these changes are better explained as a reconfiguration of linguistic priors rather than as genuine causal grounding based on visual evidence. From this perspective, future research should more fundamentally reconsider modality alignment, particularly in terms of how CF reasoning can be consistently grounded in the physical environment. Furthermore, moving beyond simple linguistic supervision, there is a need for learning paradigms that can expand the model’s reasoning capacity itself. This raises new research questions regarding how to align language-based reasoning with physical world models, and how to overcome exploration constraints imposed by pretrained priors.
% \begin{itemize}
%     % TODO : CF reasoning 의 한계에서 기인. causal confusion 관련 언급 
%     % 사실상 어떻게 world model 이 causal confusion 을 해결할 수 있을지 정당화해야함 
%     \item \textbf{From Language–Vision Alignment to Interaction-Aware Grounding.} 
%     % 본 연구의 분석은 CF reasoning이 실제 시각적 evidence와의 정합성보다는, 언어적 수준에서의 분포 재구성에 의해 작동할 가능성을 시사한다. 이는 현재의 CF reasoning이 주로 정적인 language–vision 대응에 기반하여 학습되며, 물리적 환경에서의 상태 변화나 객체 간 동적 상호작용을 충분히 반영하지 못한다는 한계로 이어진다. 이러한 관점에서, 향후 연구에서는 modality를 단순한 language–vision 정렬을 넘어, 환경과의 interaction을 포함하는 형태로 확장할 필요가 있다. 특히 world model과의 결합을 통해, 단순한 서술적 counterfactual을 넘어 실제 환경에서의 상태 전이와 그 결과를 시뮬레이션할 수 있는 방향이 요구된다. 이는 CF reasoning이 정적인 텍스트 기반 추론을 넘어, 물리적 일관성과 상호작용을 포함하는 보다 grounded한 causal reasoning으로 이어질 수 있는지에 대한 중요한 질문을 제기한다.
%     CF reasoning primarily operates through linguistic distributional reconfiguration rather than visual grounding, largely due to its reliance on static language–vision correspondences that fail to capture dynamic interactions. This limitation highlights the need to move beyond static alignment toward interaction-aware multimodal frameworks, raising the question of whether CF reasoning can ultimately achieve genuinely grounded causal reasoning.

%     % TODO : CF reasoning 의 가능성으로부터 기인.
%     \item \textbf{counterfactual reasoning with reinforcement learning.} 
%     % 본 연구에서 관찰된 바와 같이, CF reasoning은 모델의 불확실성과 탐색 특성에 변화를 유도하며, 이는 행동 분포의 재구성으로 나타날 수 있다. 이러한 특성은 CF reasoning이 reinforcement learning과 결합될 경우, 다양한 반사실적 시나리오를 탐색하는 방향으로 확장될 가능성을 시사한다. 그러나 동시에, 이러한 탐색이 pretrained prior에 의해 강하게 제약될 수 있다는 점 역시 중요한 한계로 남는다. 특히, 언어 생성의 vast action space에서는 prior 밖의 샘플이 낮은 보상을 받을 가능성이 높기 때문에, 학습 과정이 기존 분포를 강화하는 방향으로 수렴할 수 있다. 따라서 향후 연구에서는 counterfactual 기반 보상 설계나 시나리오 생성이 실제로 prior를 넘어서는 새로운 reasoning 패턴을 유도할 수 있는지, 혹은 기존 분포를 재구성하는 데 그치는지에 대한 보다 근본적인 질문이 요구된다.
%     CF reasoning reshapes action distributions by altering uncertainty and exploration characteristics, but its effects remain fundamentally constrained by pretrained priors. In the vast action space of language generation, training tends to reinforce existing distributions, raising the open question of whether counterfactual-based methods can truly induce reasoning patterns beyond the prior or merely redistribute the existing policy.
    
% \end{itemize}

\clearpage
\section{Dataset}
\label{appendix:dataset}

\subsection{Overview of the Dataset}

\begin{table*}[h]
\centering
\resizebox{\textwidth}{!}{%
\begin{tabular}{@{}l c c l r c c c c@{}}
\toprule
\textbf{Dataset} & \textbf{Year} & \textbf{Type} & \textbf{Perspect.} & \textbf{\#Frames} & \textbf{Acc. Recog.} & \textbf{Desc.} & \textbf{CF QA} & \textbf{Decision} \\
\midrule
% AIAD  & 2018 & real & variable & 10,480 & \cmark & \xmark & \xmark & \xmark \\
% CADP  & 2018 & real & roadside & 75,030* & \cmark & \xmark & \xmark & \xmark \\
VIENA$^2$ \cite{dataset:VIENA}  & 2019 & synth & vehicle & 2,250,000 & \cmark & \xmark & \xmark & \xmark \\
GTACrash \cite{dataset:GTACrash} & 2019 & synth & vehicle & 751,610 & \cmark & \xmark & \xmark & \xmark \\
% CTA  & 2020 & real & vehicle & 342,495* & \cmark & \xmark & \xmark & \xmark \\
% CCTVF  & 2020 & real & roadside & 990 & \cmark & \xmark & \xmark & \xmark \\
Argus \cite{dataset:Argus}  & 2021 & real & roadside & 120,000 & \cmark & \xmark & \xmark & \xmark \\
YoutubeCrash \cite{dataset:YoutubeCrash} & 2021 & real & vehicle & 7,720 & \cmark & \xmark & \xmark & \xmark \\
IITH Road \cite{dataset:IITH_Road} & 2022 & real & roadside & 127,138 & \cmark & \xmark & \xmark & \xmark \\
TAD \cite{dataset:TAD} & 2022 & real & roadside & 24,810 & \cmark & \xmark & \xmark & \xmark \\
MP-RAD \cite{dataset:MP-RAD} & 2023 & synth & roadside & 366,600 & \cmark & \xmark & \xmark & \xmark \\
% ADM  & 2023 & real & roadside & 3,250 & \cmark & \xmark & \xmark & \xmark \\
RiskBench \cite{dataset:riskbench} & 2024 & synth & vehicle & N/A & \cmark & \xmark & \xmark & \xmark \\
DeepAccident \cite{dataset:DeepAccident} & 2024 & synth & V2V \& V2I & 57,000 & \cmark & \xmark & \xmark & \xmark \\
MM-AU \cite{dataset:mm_au} & 2024 & real & vehicle & 2,190,000 & \cmark & \cmark & \xmark & \xmark \\
TUMTraf-Accid3nD \cite{dataset:TUMTraf_Accid3nD} & 2025 & real & roadside & 111,656 & \cmark & \cmark & \xmark & \xmark \\
% \rowcolor{gray!10} % 마지막 줄 강조 (회색 배경)
\textbf{Ours} & 2026 & real & vehicle & \textbf{2,620,356} & \cmark & \cmark & \cmark & \cmark \\
\bottomrule
\end{tabular}%
}
\caption{Comparison of our testbed dataset with existing accident and traffic datasets. Unlike previous works limited to recognition, our dataset supports high-level reasoning tasks including description, counterfactual QA, and decision making.}
\label{tab:dataset_comparison}
\end{table*}
% 사고 시나리오의 경우, 정상 주행 시나리오에 비해 confounder 로 인한 라벨링의 어려움이 덜하고 counterfactuals 포착 및 평가에 더 적합함. 

The testbed dataset used in this study was constructed by integrating \cite{dataset:ccd, dataset:coool, dataset:nexar} and applying captioning to each video. Importantly, each video in these datasets has been annotated by domain experts with labels indicating whether an accident occurs. Each dataset provides predefined labels indicating whether a given video corresponds to an accident scenario or a normal driving scenario. However, these existing driving accident datasets primarily focus on recognition of the scene and do not include driving actions or counterfactual supervision. As a result, they lack counterfactual captions and decision required to address the central question of this study: whether counterfactual reasoning truly affects safety-relevant performance. 


% 총 프레임 수는 2620356, accident recog., description, counterfactual QA 여부, decision 포함 여부
\begin{table}[h]
\centering
\small
\begin{tabular}{lccc}
\toprule
\textbf{Decision Class} & \textbf{Train} & \textbf{Test} & \textbf{Total} \\
\midrule
stop & 1,026 (15.4\%) & 336 (25.0\%) & 1,362 \\
decelerate     & 1,423 (21.4\%) & 336 (25.0\%) & 1,759 \\
keep      & 4,212 (63.2\%) & 672 (50.0\%) & 4,884 \\
\midrule
\textbf{Total} & \textbf{6,661} & \textbf{1,344} & \textbf{8,005} \\
\bottomrule
\end{tabular}
\caption{Dataset composition. Test set is intentionally 
balanced for robust evaluation of safety-critical decisions.}
\end{table}


\subsection{Automated Annotation Pipeline}

\begin{figure}[h]
    \centering
    \includegraphics[width=\columnwidth]{figure/dataset/automated_pipeline.png}
    \caption{The automated annotation pipeline for the dataset.}
    \label{fig:automated_pipeline}
\end{figure}

To address these limitations, we introduce a multi-stage automated annotation pipeline to construct a dataset for analyzing whether counterfactual reasoning truly affects safety-relevant performance from existing accident driving videos. Drawing inspiration from the visual question answering methodology proposed in \cite{dataset:lingoqa}, we leverage the advanced multimodal capabilities of \texttt{Qwen2.5-VL-32B-Instruct} to generate structured captions \cite{e2e-vlm:orion, causal:cf_vla}. Our annotation process produces six key components for each video: \texttt{accident occurrence, ego-vehicle involvement, scene description, counterfactual question, counterfactual answer}, and \texttt{driving decision}. The pipeline consists of three sequential steps:

\begin{itemize}[leftmargin=*, noitemsep, topsep=0pt]
    \item \textbf{Step 1: Scene Recognition and Description.} 
    First, we feed the raw accident video into \texttt{Qwen2.5-VL-32B-Instruct} to extract fundamental scene information. The model is prompted to identify whether an accident occurred (\texttt{accident occurrence}) and whether the ego-vehicle was involved (\texttt{ego involvement}). Simultaneously, it generates a detailed textual explanation of the traffic scenario (\texttt{description}), focusing on the behavior of road agents and environmental conditions.

    \item \textbf{Step 2: Counterfactual Reasoning Generation.} 
    Based on the \texttt{description} obtained in Step 1, we guide the model to generate a \texttt{counterfactual question} and its corresponding \texttt{counterfactual answer}. This step aims to augment the dataset with reasoning capabilities by asking "what-if" scenarios (e.g., \textit{"what would have happened if the ego-vehicle had not braked?"}), thereby enhancing the model's understanding of causal relationships in safety-critical situations.

    \item \textbf{Step 3: Decision Making.} 
    Finally, we synthesize the \texttt{decision} for the original video. The model aggregates all prior information—accident status, ego-involvement, scene description, and counterfactual reasoning—to deduce the optimal driving maneuver (e.g., stop, decelerate, keep). This ensures that the final decision is logically grounded in both the observed reality and potential alternative outcomes.
\end{itemize}


% TODO : critical
\subsection{Quality Assurance}

\begin{figure}[h]
    \centering
    \includegraphics[width=\columnwidth]{figure/dataset/human_evaluation_tool.png}
    \caption{Diagnosis tool for quality assessment.}
    \label{fig:human_eval}
\end{figure}

To ensure the reliability of our dataset while maintaining the scalability of the construction process, we implemented a rigorous quality assurance protocol. Unlike traditional datasets that rely on labor-intensive manual labeling, we leverage the zero-shot reasoning capabilities of \texttt{Qwen2.5-VL-32B-Instruct} to establish an automated pipeline, which we then statistically validate through a detailed human audit.

\paragraph{Verification Protocol.} 
We randomly sampled 100 videos from our dataset for a formal verification. Independent annotators evaluated each sample in terms of factual correctness (i.e., the accuracy of the scene description relative to the visual content), counterfactual plausibility (i.e., the logical and physical soundness of the hypothetical scenarios), and decision appropriateness (i.e., the safety and optimality of the predicted driving actions).

\paragraph{Quantitative Validation.} 
%As summarized in Table~\ref{tab:human_eval_results}, 
The automated pipeline achieved an 86.0\% usability rate. The remaining 14.0\% exhibited major errors or hallucinations. While we acknowledge these imperfections, an 86\% reliability rate in complex, multi-modal reasoning tasks is highly competitive and justifies the use of raw model-generated data for large-scale training and benchmarking \cite{dataset:lingoqa}.

\begin{table*}[h]
\centering
\begin{tabular}{lcc}
\toprule
\textbf{Category} & \textbf{Count} & \textbf{Percentage (\%)} \\
\midrule
Correct & 77 & 77.0\% \\
Minor Errors (e.g., phrasing) & 9 & 9.0\% \\
Major Errors (e.g., logic flaws) & 8 & 8.0\% \\
Hallucinations (e.g., phantom objects) & 6 & 6.0\% \\
\midrule
\textbf{Usable (Correct + Minor)} & \textbf{86} & \textbf{86.0\%} \\
\bottomrule
\end{tabular}
\caption{Human audit results ($N=100$). Samples categorized as 'Correct' or 'Minor Error' are considered usable for robust training and evaluation.}
\label{tab:human_eval_results}
\end{table*}

\begin{table*}[h]
\centering
\small
\resizebox{\textwidth}{!}{
\begin{tabular}{l c | cc cc ccc}
\toprule
\multirow{2}{*}{\textbf{Model}} & \multirow{2}{*}{\textbf{Metric}} & \multicolumn{2}{c}{\textbf{Accident Recog.}} & \multicolumn{2}{c}{\textbf{Ego Inclusion}} & \multicolumn{3}{c}{\textbf{Decision}} \\
 & & Yes & No & Involv. & Not In. & Stop & Decel. & Keep \\
\midrule

\multirow{3}{*}{\textbf{InternVL2-1B}}
 & P@1 w/o & 69.4 & 58.0 & 58.0 & 69.4 & 33.3 & 32.1 & 71.4 \\
 & P@1 w/  & 66.7 & 30.0 & 66.7 & 30.0 & 22.2 & 32.1 & 59.2 \\
 & \cellcolor{gray!10} $\Delta$ & -2.7 & -28.0 & +8.7 & -39.4 & -11.1 & 0.0 & -12.2 \\
\midrule

\multirow{3}{*}{\textbf{InternVL2-2B}}
 & P@1 w/o & 41.7 & 66.6 & 41.7 & 66.6 & 11.1 & 25.0 & 83.7 \\
 & P@1 w/  & 50.0 & 56.0 & 50.0 & 56.0 & 22.2 & 32.1 & 83.7 \\
 & \cellcolor{gray!10} $\Delta$ & +8.3 & -10.6 & +8.3 & -10.6 & +11.1 & +7.1 & 0.0 \\
\midrule

\multirow{3}{*}{\textbf{InternVL2.5-1B}}
 & P@1 w/o & 44.4 & 60.0 & 44.4 & 60.0 & 22.2 & 28.6 & 69.4 \\
 & P@1 w/  & 52.8 & 60.0 & 52.8 & 60.0 & 44.4 & 25.0 & 73.5 \\
 & \cellcolor{gray!10} $\Delta$ & +8.4 & 0.0 & +8.4 & 0.0 & +22.2 & -3.6 & +4.1 \\
\midrule

\multirow{3}{*}{\textbf{Qwen2.5-VL-3B}}
 & P@1 w/o & 63.9 & 42.0 & 55.6 & 44.0 & 33.3 & 46.4 & 55.1 \\
 & P@1 w/  & 75.0 & 38.0 & 75.0 & 38.0 & 55.6 & 21.4 & 44.9 \\
 & \cellcolor{gray!10} $\Delta$ & +11.1 & -4.0 & +19.4 & -6.0 & +22.3 & -25.0 & -10.2 \\
\midrule

\multirow{3}{*}{\textbf{VideoLLaMA2-7B}}
 & P@1 w/o & 47.2 & 68.0 & 47.2 & 68.0 & 11.1 & 17.9 & 77.6 \\
 & P@1 w/  & 77.8 & 34.0 & 75.0 & 34.0 & 33.3 & 28.6 & 46.9 \\
 & \cellcolor{gray!10} $\Delta$ & +30.6 & -34.0 & +27.8 & -34.0 & +22.2 & +10.7 & -30.7 \\
\midrule

\end{tabular}
}
\caption{Pass@1 comparison w/ and w/o counterfactual in SFT for random samples.}
\label{tab:passk_sft_quality_validation}
\end{table*}


To verify robustness to annotation noise, we re-evaluated on the 86 high-quality samples confirmed through human audit. The Pass@1 results demonstrate identical directional trends to those observed in the full dataset, confirming that the 14\% noise does not systematically bias our conclusions.

\clearpage
\subsection{Statistical Analysis of the Dataset}


\begin{figure}[h]
    \centering
    \includegraphics[width=\columnwidth]{figure/dataset/figure1_basic_distribution.pdf}
    \caption{\textbf{Overview of dataset statistics and label distributions.} The plots illustrate the imbalance in accident occurrence (top-left) and ego-vehicle involvement (bottom-left). The horizontal bar chart (right) compares the distribution of driving decisions (stop, decelerate, keep) across the training (blue) and testing (red) splits, demonstrating a consistent class distribution between the two sets.}
    \label{fig:dataset_static_fig1_clss}
\end{figure}

\begin{figure}[h]
    \centering
    \includegraphics[width=\columnwidth]{figure/dataset/figure2_decision_comparison.pdf}
    \caption{\textbf{Conditional distribution of driving decisions with respect to accident occurrence and ego-vehicle involvement.} The heatmaps display column-normalized percentages to analyze the correlation between decisions and scenario attributes. The top row shows the relationship between decisions and accidents, while the bottom row depicts ego-involvement. Left columns represent the training set (blue scale), and right columns represent the test set (red scale).}
    \label{fig:dataset_static_fig2_cmap}
\end{figure}

\begin{figure}[h]
    \centering
    \includegraphics[width=\columnwidth]{figure/dataset/figure3_word_count_metrics.pdf}
    \caption{\textbf{Analysis of linguistic characteristics and token distributions.} (top-left) Kernel Density Estimation (KDE) of total word counts showing alignment between training and testing splits. (top-right) Average word counts broken down by descriptions, counterfactual Q\&A. (bottom) Violin plots illustrating the distribution of word counts for counterfactual answers (bottom-left) and questions (bottom-right). Blue indicates the training set, and red indicates the test set.}
    \label{fig:dataset_static_fig3_word.dist}
\end{figure}

\begin{figure}[h]
    \centering
    \includegraphics[width=\columnwidth]{figure/dataset/figure4_scenario_analysis.pdf}
    \caption{\textbf{Distribution of primary accident causes.} The horizontal bar charts display the frequency of the top annotated accident causes (e.g., control loss, perception failure) for the training (left, blue) and testing (right, red) sets. The semantic distribution of accident scenarios is preserved across the splits.}
    \label{fig:dataset_static_fig4_scene}
\end{figure}


\clearpage 
\section{Prompt Templates}
\label{appendix:prompts}
% \paragraph{Prompts for Automated Captioning}
\begin{tcolorbox}[title=description and counterfactual Q\&A generation for normal video, breakable]
\small
\textbf{System prompt:}
\begin{verbatim}

You are given a video of driving scene with no accident. 
For this video, **there is no accident**. 
Describe the following information in natural language sentences 
(each item should be 10-20 words):

1. Accident occurrence (always 0)
2. Whether the ego-car is involved (always 0)
3. Scene description 
(describe what is actually happening in the video)
4. Counterfactual question 
(what-if scenario)
5. Counterfactual answer 
(what would happen in the what-if scenario)

Use complete sentences to report all six items.

Video Description:
<scene summary or keyframe captions here>
Ego-car involvement: <yes/no>

Strictly follow this output format: 
1, 1, <accident>...</accident>, <accident>...</accident>, 
<no_accident>...</no_accident>, <no_accident>...</no_accident>

For counterfactual question and answer, 
create diverse what-if scenarios beyond vehicle swerving. 
Consider weather, pedestrians, traffic signals, road conditions, 
or mechanical failures.

Now, generate the annotation for the given video 
as a single comma-separated line. 
Do not add any extra text.
\end{verbatim}
\end{tcolorbox}


\begin{tcolorbox}[title=description and counterfactual Q\&A generation  for accident video, breakable]
\small
\textbf{System prompt:}
\begin{verbatim}
You are given a video of an accident. 
For this video, **an accident occurs or harzard scene exists**. 
Describe the following information in natural language sentences 
(each item should be 10-20 words):

1. Accident occurrence (always 1)
2. Whether the ego-car is involved (0 for no, 1 for yes)
3. Accident result
4. Accident cause
5. Counterfactual question (what-if scenario)
6. Counterfactual answer 
(what would happen if no accident occurred)

Use complete sentences to report all six items.

Video Description:
<scene summary or keyframe captions here>
Ego-car involvement: <yes/no>

Strictly follow this output format:
Output: 1, (0 for no, 1 for yes), 
<accident>...</accident>, <accident>...</accident>, 
<no_accident>...</no_accident>, <no_accident>...</no_accident>

Now, generate the annotation for the given video 
as a single comma-separated line. 
Do not add any extra text.
\end{verbatim}
\end{tcolorbox}

\begin{tcolorbox}[title=decision generation, breakable]
\small
\textbf{System prompt:}
\begin{verbatim}
You are an expert driving policy annotator. 
You are given a video and its captions describing a driving scene. 
The video may show either normal driving or the moments just 
before a potential accident. 

The caption always follows this format:

1. Accident occurrence (0 for no accident, 1 for accident)
2. Whether the ego-car is involved (0 for no, 1 for yes)
3. Accident result
4. Accident cause
5. Counterfactual question (what-if scenario)
6. Counterfactual answer 
(what would happen if no accident occurred)

Examples:

# Accident occurred
Accident occurrence: 1
Ego involved: 0 or 1
Accident result: 
    a car in front of ego-car collides with a pedestrian
Accident cause: 
    the pedestrian violates traffic signal
Counterfactual question: 
    what if the pedestrian does not violate the signal?
Counterfactual answer: 
    there would be no collision

# No accident
Accident occurrence: 0
Ego involved: 0
Accident result: 
    no accident
Accident cause: 
    the car in front of ego-car turns right
Counterfactual question: 
    what if the car suddenly stopped?
Counterfactual answer: 
    there would be a collision

---

Caption for the given video:
{data['conversations'][1]['value']}

Your task:
Return the **driving decision** as a single action class.  
The possible action classes are:

- STOP
- DECELERATE
- KEEP

Rules:
1. If an accident occurs and the ego-car is involved: 
return the action the ego-car should have taken before the accident 
2. If an accident occurs but the ego-car is not involved: 
return the safe action to avoid the accident 
3. If no accident occurs or no action is required: 
return <decision>KEEP</decision>.
4. You may refer to the counterfactual question and answer 
if it helps decide the correct action.

Output format:
<decision>ACTION</decision>

Return only a single line in this format. 
Do not add any explanation or extra text.
\end{verbatim}
\end{tcolorbox}


% \subsection{Prompt for Experiment}

\begin{tcolorbox}[title=Fine-Tuning and Test-Time Prompts for model w/o CF reasoning, breakable]
\small
\textbf{System prompt:}
\begin{verbatim}


You are given a video of a driving scene. 
The video may show normal driving 
or the moments just before a potential accident. 

Describe the following information in natural language sentences 
(each item should be 10-20 words):
1. Accident occurrence (0 for no accident, 1 for accident)
2. Whether the ego-car is involved (0 for no, 1 for yes)
3. Accident result
4. Accident cause
5. Decision for given scene

Use complete sentences to report all six items.

Video Description:
<scene summary or keyframe captions here>

Examples:

# Accident occurred
Accident occurrence: 1
Ego involved: 0 or 1
Accident result: 
    a car in front of ego-car collides with a pedestrian
Accident cause: 
    the pedestrian violates traffic signal

# No accident
Accident occurrence: 0
Ego involved: 0
Accident result: 
    no accident
Accident cause: 
    the car in front of ego-car turns right
Now, describe the following video in the same sentence-style format:
<scene summary or keyframe captions of your target video>

Format of Accident occurred:
Output: 1, (0 for no, 1 for yes), 
<accident>Accident result</accident>, 
<accident>Accident cause</accident>, 
<decision>DECISION</decision>

Format of No accident:
Output: 0, 0, 
<no_accident>Accident result</no_accident>, 
<no_accident>Accident cause</no_accident>, 
<decision>DECISION</decision>

Now, generate the response for the given video 
as a **single comma-separated line**. 
Do not add any extra text.
\end{verbatim}
\end{tcolorbox}


\begin{tcolorbox}[title=Fine-Tuning and Test-Time Prompts for model w/ CF reasoning, breakable]
\small
\textbf{System prompt:}
\begin{verbatim}


You are given a video of a driving scene. 
The video may show normal driving 
or the moments just before a potential accident. 

Describe the following information in natural language sentences 
(each item should be 10-20 words):
1. Accident occurrence (0 for no accident, 1 for accident)
2. Whether the ego-car is involved (0 for no, 1 for yes)
3. Accident result
4. Accident cause
5. Counterfactual question 
(what-if scenario)
6. Counterfactual answer 
(what would happen if no accident occurred)
7. Decision for given scene

Use complete sentences to report all six items.

Video Description:
<scene summary or keyframe captions here>

Examples:

# Accident occurred
Accident occurrence: 1
Ego involved: 0 or 1
Accident result: 
    a car in front of ego-car collides with a pedestrian
Accident cause: 
    the pedestrian violates traffic signal
Counterfactual question: 
    what if the pedestrian does not violate the signal?
Counterfactual answer: 
    there would be no collision

# No accident
Accident occurrence: 0
Ego involved: 0
Accident result: 
    no accident
Accident cause: 
    the car in front of ego-car turns right
Counterfactual question: 
    What if the car suddenly stopped?
Counterfactual answer: 
    there would be a collision
Now, describe the following video in the same sentence-style format:
<scene summary or keyframe captions of your target video>

Format of Accident occurred:
Output: 1, (0 for no, 1 for yes), 
<accident>Accident result</accident>, 
<accident>Accident cause</accident>, 
<no_accident>Counterfactual question</no_accident>, 
<no_accident>Counterfactual Answer</no_accident>
<decision>DECISION</decision>

Format of No accident:
Output: 0, 0, 
<no_accident>Accident result</no_accident>, 
<no_accident>Accident cause</no_accident>, 
<accident>Counterfactual Question</accident>, 
<accident>Counterfactual Answer</accident>
<decision>DECISION</decision>

Now, generate the response for the given video 
as a **single comma-separated line**. 
Do not add any extra text.
\end{verbatim}
\end{tcolorbox}



\begin{tcolorbox}[title=Token Activation Map Analysis Prompts, breakable]
\small
\textbf{System prompt:}
\begin{verbatim}
# perceptual prompt
What is the most critical object the ego-car should be cautious of 
in this driving scene? 

# decision prompt 
Based on the current visual scene, what is the immediate action 
required for the ego-car?

# counterfactual prompt
(counterfactual question on the main factor based on the GT scene description)
(e.g.) What if the car suddenly stopped?
\end{verbatim}
\end{tcolorbox}

\clearpage
\section{Reward Functions}
\label{appendix:reward_functions}
To align the model's outputs with safety-critical driving standards, we design a comprehensive reward function $R$ composed of six sub-rewards. Let $y$ denote the generated completion and $y^*$ denote the ground truth. The total reward is defined as a weighted sum:

\begin{equation}
    R = w_{\text{acc}} R_{\text{acc}} + w_{\text{ego}} R_{\text{ego}} + w_{\text{dec}} R_{\text{dec}} + w_{\text{safe}} R_{\text{safe}} + w_{\text{corr}} R_{\text{corr}} + w_{\text{fmt}} R_{\text{fmt}}
\label{eq:total_reward}
\end{equation}

\noindent where the weights are set to $w_{\text{acc}}=0.1$, $w_{\text{ego}}=0.1$, $w_{\text{dec}}=0.4$, $w_{\text{safe}}=0.2$, $w_{\text{corr}}=0.1$, and $w_{\text{fmt}}=0.1$, prioritizing the correct driving decision and safety.

\paragraph{Recognition Rewards ($R_{\text{acc}}, R_{\text{ego}}$)}
We verify the fundamental perception capabilities: accident occurrence ($y_{\text{acc}}$) and ego-vehicle involvement ($y_{\text{ego}}$). These rewards utilize the indicator function $\mathbb{I}(\cdot)$, which equals 1 if the condition is met and 0 otherwise:
\begin{equation}
    R_{\text{acc}} = \mathbb{I}(y_{\text{acc}} = y^*_{\text{acc}}), \quad R_{\text{ego}} = \mathbb{I}(y_{\text{ego}} = y^*_{\text{ego}})
\end{equation}

\paragraph{Decision Reward ($R_{\text{dec}}$)}
The core driving decision (e.g., decelerate, stop) is extracted from the \texttt{<decision>} tags.
\begin{equation}
    R_{\text{dec}} = \mathbb{I}(y_{\text{dec}} = y^*_{\text{dec}})
\end{equation}


\paragraph{Safety Critical Reward ($R_{\text{safe}}$)}
To penalize dangerous failures, we introduce a specialized reward for ``stop'' scenarios. If the ground truth requires a stop ($y^*_{\text{dec}} = \text{stop}$), the model receives a penalty for missing it:
\begin{equation}
    R_{\text{safe}} = 
    \begin{cases} 
    1.0 & \text{if } y^*_{\text{dec}} = \text{stop} \land y_{\text{dec}} = \text{stop} \\
    -0.5 & \text{if } y^*_{\text{dec}} = \text{stop} \land y_{\text{dec}} \neq \text{stop} \\
    0.0 & \text{otherwise}
    \end{cases}
\end{equation}

% TODO : 좀 더 구체화 
\paragraph{Consistency Reward ($R_{\text{corr}}$)}
To ensure logical consistency between the classification flag and the textual description, $R_{\text{corr}}$ verifies whether the predicted accident flag $y_{\text{acc}}$ matches the ground truth and aligns with the corresponding XML tags. Let $T(y)$ be the set of tags generated in the description. The reward is defined as:

\begin{equation}
    R_{\text{corr}} = \mathbb{I}(y_{\text{acc}} = y^*_{\text{acc}}) \cdot \mathbb{I}\left( 
    \begin{aligned} 
    &(y_{\text{acc}} = 1 \land \texttt{<accident>} \in T(y)) \lor \\ 
    &(y_{\text{acc}} = 0 \land \texttt{<no\_accident>} \in T(y)) 
    \end{aligned} 
    \right)
\end{equation}

\paragraph{Format Rewards ($R_{\text{fmt}}$)} 
This reward ensures structural validity by verifying the number of parsed segments $|S_y|$. The expected length differs depending on the reasoning mode. For models utilizing counterfactual reasoning, we require extended segments for the traces:

\begin{equation}
    R_{\text{fmt}} = 
    \begin{cases} 
    \mathbb{I}(|S_y| = 7) & \text{if w/ CF reasoning)} \\
    \mathbb{I}(|S_y| = 5) & \text{otherwise}
    \end{cases}
\end{equation}

% TODO : safety_critical_reward 가 confound 로 작용할 수 있지만, 치명타는 아님. 달리 얘기하면, safety_reward 가 confound 로 작용하여 safety-bias 를 유발하는 요인으로 오해될 수 있다는 뜻. 그러나 Qwen2.5-VL-3B w/o CF 를 기준으로 보면, RFT 이후에 오히려 stop 에 대해서 pass@k 성능이 감소함
The $R_{\text{safe}}$ assigns additional weight to emergency stop predictions, which may amplify the safety prior observed in SFT. However, since both w/ CF and w/o CF reference models are trained under identical reward structures, the observed difference in RFT action distribution shift is attributable to the CF reasoning effect in the reference model rather than the reward design alone. The persistence of the SFT-RFT gap suggests that CF reasoning establishes a prior that interacts with, rather than is replaced by, the reward signal.

\paragraph{Disentangling CF Reasoning from Reward Design.} To verify whether the observed action distribution shift is attributable to the reward design rather than CF reasoning itself, we conduct an ablation study by removing the safety-critical reward $R_{\text{safe}}$ from the reward function. We select Qwen2.5-VL-3B (GRPO) as the target model, as it exhibited the strongest CF-induced bias in Table~\ref{tab:passk_rft}. To compensate for the removed $R_{\text{safe}}$, we increase the decision reward weight from $w_{\text{dec}}=0.4$ to $w_{\text{dec}}=0.6$, ensuring that the total reward scale remains comparable to the original setting.

\begin{figure}[h]
    \centering
    \includegraphics[width=\columnwidth]{figure/passk/task_pass_at_k_grid_Rsafe-ablation.pdf}
    \caption{Class-wise pass@$k$ performance across model scales and CF reasoning in ablation on $R_{\text{safe}}$.}
    \label{fig:passk_Rsafe}
\end{figure}

\begin{table*}[h]
\centering
\begin{tabular}{c | cc cc ccc}
\toprule
\multirow{2}{*}{Metric} & \multicolumn{2}{c}{\textbf{Accident Recog.}} & \multicolumn{2}{c}{\textbf{Ego Inclusion}} & \multicolumn{3}{c}{\textbf{Decision}} \\
 & Yes & No & Involv. & Not In. & Stop & Decel. & Keep \\
\midrule
P@1 w/o CF & 50.6 & 57.3 & 46.7 & 60.7 & 19.3 & 32.3 & 65.8 \\
P@1 w/ CF  & 64.2 & 44.8 & 61.5 & 46.0 & 37.0 & 28.5 & 53.2 \\
\cellcolor{gray!10} $\Delta$ & +13.6 & -12.5 & +14.8 & -14.7 & +17.7 & -3.8 & -12.6 \\
\cmidrule{1-8}
P@16 w/o CF & 95.7 & 93.9 & 95.7 & 97.9 & 69.5 & 91.8 & 96.9 \\
P@16 w/ CF  & 99.7 & 97.6 & 99.7 & 98.7 & 94.5 & 92.6 & 98.3 \\
\cellcolor{gray!10} $\Delta$ & +4.0 & +3.7 & +4.0 & +0.8 & +25.0 & +0.8 & +1.4 \\
\bottomrule
\end{tabular}
\caption{
Effect of counterfactual (CF) reasoning \textbf{without $R_{\text{safe}}$} across different $k$ on Qwen2.5-VL-3B (GRPO). While the magnitude of the shift is reduced compared to the full setting, consistent improvements in accident-related and safety-critical decisions (e.g., Stop) persist across both P@1 and P@16, indicating that the effect is not limited to low-$k$ sampling.
}
\label{tab:rsafe_ablation}
\end{table*}


As shown in Table~\ref{tab:rsafe_ablation} and Figure ~\ref{fig:passk_Rsafe}, the distribution shift induced by CF reasoning persists even in the absence of $R_{\text{safe}}$: the Stop decision rate increases by $+17.7$\%p under the CF condition compared to the non-CF condition.
While the magnitude of the shift is reduced compared to the full reward setting ($+31.0$\%p), the direction remains consistent, indicating that $R_{\text{safe}}$ amplifies the bias but does not cause it. These results provide direct evidence that the action distribution shift originates from CF reasoning itself, rather than from the reward design.




\clearpage 
\section{Model Selection Criteria}
\label{appendix:model_selection_criteria}


To ensure the generality of our findings, we select backbone models according to three principles:
(1) prevalence in recent Vision-Language-Action (VLA) literature,
(2) architectural diversity in multimodal fusion strategies, and
(3) scalability across different parameter regimes under practical computational constraints.
% FIXME:
% 사실상 \cite{e2e-vlm:are_vlms_ready_for_ad} 참고. 여기서도 model size 10B 이하만 선택됨요

\paragraph{Architectural Coverage.}
We adopt models from three representative multimodal families that are widely used in contemporary autonomous driving research:

\begin{itemize}
    \item \textbf{InternVL2-1B, InternVL2-2B and InternVL2.5-1B:}
    These models have recently emerged as strong open-source backbones in driving-related VLA systems. 
    Using multiple scales within the same architectural family enables controlled analysis of parameter scaling effects.
    
    \item \textbf{Qwen2.5-VL-3B:}
    A widely adopted backbone in recent large-scale VLA frameworks, representing a distinct multimodal alignment strategy compared to InternVL.
    
    \item \textbf{VideoLLaMA2-7B:}
    A temporally-aware video-language model, included to evaluate whether the observed trends persist in video-native architectures with larger parameter scales.
\end{itemize}

Collectively, our selection spans diverse multimodal fusion designs and parameter regimes ranging from 1B to 7B parameters, mitigating the risk that the observed effects are backbone-specific artifacts.

\begin{table*}[h]
    \centering
    \footnotesize 
    \renewcommand{\arraystretch}{1.3} % 행 간격은 유지 (가독성)
    \setlength{\tabcolsep}{3pt} 
    
    \begin{tabularx}{\textwidth}{ 
        >{\raggedright\arraybackslash}p{2.2cm} 
        >{\raggedright\arraybackslash}p{4.5cm} 
        >{\raggedright\arraybackslash}X 
    }
    \toprule
    \textbf{Category} & \textbf{Rationale for Selection} & \textbf{Related Works} \\
    \midrule

    % --- LLaMA & Mistral ---
    \textbf{LLaMA \& Mistral} 
    & Standard baseline in early VLA research. We adopt \textbf{VideoLLaMA2-7B} to represent this lineage.
    & DriveMLM \cite{e2e-vlm:drivemlm}, RDA-Driver \cite{e2e-vlm:rda-driver}, DriveLM \cite{e2e-vlm:drivelm}, DriveGPT4 \cite{e2e-vlm:drivegpt4}, VLAAD \cite{e2e-vlm:vlaad}, OccLLaMA \cite{e2e-vlm:occllama}, OmniDrive \cite{causal:e2e-vlm:omnidrive}, LMDrive \cite{e2e-vlm:lmdrive}, BEVDriver \cite{e2e-vlm:bevdriver}, CoVLA-Agent \cite{e2e-vlm:covla-agent}, DSDrive \cite{e2e-vlm:dsdrive} \\

    \midrule

    % --- Qwen-VL ---
    \textbf{Qwen Family} 
    & Dominant in 2025 research. We adopt \textbf{Qwen2.5-VL-3B} to align with SOTA trends.
    & ImpromptuVLA \cite{e2e-vlm:impromptuvla}, DriveAgent-R1 \cite{e2e-vlm:driveagent-r1}, FastDriveVLA \cite{e2e-vlm:fastdrivevla}, AutoDrive-R$^2$ \cite{e2e-vlm:autodrive-r2}, CoReVLA \cite{e2e-vlm:corevla}, AutoVLA \cite{e2e-vlm:autovla}, E3AD \cite{e2e-vlm:e3ad}, UniUGP \cite{e2e-vlm:uniugp}, Reasoning-VLA \cite{e2e-vlm:reasoning-vla} \\

    \midrule

    % --- InternVL ---
    \textbf{InternVL Family} 
    & Vision-Centric alternative. We adopt \textbf{InternVL2} to verify scalability.
    & Drive-R1 \cite{e2e-vlm:drive-r1}, SimLingo \cite{causal:e2e-vlm:simlingo}, VDRive \cite{e2e-vlm:vdrive}, DriVLM \cite{e2e-vlm:drivlm}, Li et al. \cite{e2e-vlm:driving_with_internvl_challange} \\

    \bottomrule
    \end{tabularx}
    \caption{Summary of VLA Models and Related Works. We categorize models based on their architectural lineage.}
    \label{tab:backbone_summary}
\end{table*}

\paragraph{Use of Public Resources.}
We acknowledge the use of the following publicly available resources in this work. :

\begin{itemize}

    \item InternVL2-1B \footnote{\url{https://huggingface.co/OpenGVLab/InternVL2-1B}} \dotfill Apache License 2.0

    \item InternVL2-2B \footnote{\url{https://huggingface.co/OpenGVLab/InternVL2-2B}} \dotfill Apache License 2.0

    \item InternVL2.5-1B \footnote{\url{https://huggingface.co/OpenGVLab/InternVL2_5-1B}} \dotfill Apache License 2.0

    \item Qwen2.5-VL-3B \footnote{\url{https://huggingface.co/Qwen/Qwen2.5-VL-3B-Instruct}} \dotfill Apache License 2.0

    \item VideoLLaMA2-7B \footnote{\url{https://huggingface.co/DAMO-NLP-SG/VideoLLaMA2-7B-Base}} \dotfill Apache License 2.0

\end{itemize}


\clearpage
\section{Training Details}
\label{appendix:hyperparams}

For supervised fine-tuning (SFT), we utilized four NVIDIA A5000 GPUs, whereas reinforcement fine-tuning (RFT) was conducted using a single NVIDIA L40S GPU.

\begin{table*}[h]
\centering
\begin{tabular}{llll}
\toprule
\textbf{Setting} & \textbf{Model} & \textbf{Hyperparameter} & \textbf{Value} \\
\midrule

\multirow{30}{*}{SFT}
& \multirow{6}{*}{InternVL2-1B}
& Batch size & 4 \\
& & Precision & \texttt{bf16} \\
& & Epoch & 5 \\
& & Learning rate & $1\times10^{-5}$ \\
& & Weight decay & 0.05 \\
& & \# of frames & 8 \\ 
\cmidrule(lr){2-4}

& \multirow{6}{*}{InternVL2-2B}
& Batch size & 4 \\
& & Precision & \texttt{bf16} \\
& & Epoch & 5 \\
& & Learning rate & $1\times10^{-5}$ \\
& & Weight decay & 0.05 \\
& & \# of frames & 8 \\ 
\cmidrule(lr){2-4}

& \multirow{6}{*}{InternVL2.5-1B}
& Batch size & 4 \\
& & Precision & \texttt{bf16} \\
& & Epoch & 5 \\
& & Learning rate & $1\times10^{-5}$ \\
& & Weight decay & 0.05 \\
& & \# of frames & 8 \\ 
\cmidrule(lr){2-4}

& \multirow{6}{*}{Qwen2.5-VL-3B}
& Batch size & 2 \\
& & Precision & \texttt{bf16} \\
& & Epoch & 5 \\
& & Learning rate & $1\times10^{-5}$ \\
& & Weight decay & 0.01 \\
& & \# of frames & 8 \\ 
\cmidrule(lr){2-4}

& \multirow{6}{*}{VideoLLaMA2-7B}
& Batch size & 2 \\
& & Precision & \texttt{bf16} \\
& & Epoch & 5 \\
& & Learning rate & $2\times10^{-5}$ \\
& & Weight decay & 0.01 \\
& & \# of frames & 8 \\ 

\midrule

\multirow{10}{*}{RFT}
& \multirow{5}{*}{InternVL2.5-1B (GRPO)}
& Batch size & 2 \\
& & Temperature & 0.9 \\ 
& & KL loss coefficient & $4\times10^{-2}$ \\
& & Group size & 2 \\
& & Learning rate & $1\times10^{-6}$ \\
\cmidrule(lr){2-4}

& \multirow{5}{*}{Qwen2.5-VL-3B (GRPO, GDPO)}
& Batch size & 2 \\
& & Temperature & 1.2 \\
& & KL loss coefficient & $4\times10^{-2}$ \\
& & Group size & 2 \\
& & Learning rate & $1\times10^{-6}$ \\

\bottomrule
\end{tabular}
\caption{Training hyperparameters for SFT and RFT.}
\label{tab:hyperparams_all}
\end{table*}



\clearpage
\section{Pseudo Codes}
\label{appendix:pseudo_codes}

\subsection{Policy Entropy}
\begin{lstlisting}[label={lst:policy_entropy}]
def calculate_policy_entropy(logits, completion_mask, temperature=1.0):
    """
    Computes the length-normalized policy entropy for generated completions.
    
    Args:
        logits (Tensor): Shape (Batch, Seq_Len, Vocab_Size)
        completion_mask (Tensor): Shape (Batch, Seq_Len). 1 for completion, 0 for prompt.
        temperature (float): Scaling factor for logits.
    
    Returns:
        float: Average entropy per token (invariant to sequence length).
    """
    
    # 1. Temperature Scaling
    # Adjusts the sharpness of the distribution.
    logits = logits / temperature

    # 2. Compute Entropy Per Token
    # H(t) = - sum(p(x) * log p(x))
    probs = softmax(logits, dim=-1)
    log_probs = log_softmax(logits, dim=-1)
    
    # Shape: (Batch, Seq_Len)
    per_token_entropy = -sum(probs * log_probs, dim=-1)

    # 3. Apply Completion Mask
    # Filter out prompt tokens (entropy = 0 for prompt positions)
    masked_entropy = per_token_entropy * completion_mask

    # 4. Length Normalization (Critical Step)
    # Calculate effective length of the generated completion
    sequence_length = sum(completion_mask, dim=-1)
    
    # Divide total entropy by sequence length
    # This ensures the metric reflects 'uncertainty density' rather than total uncertainty.
    policy_entropy = sum(masked_entropy, dim=-1) / clamp(sequence_length, min=1.0)

    return mean(policy_entropy) # Average over batch
\end{lstlisting}

\subsection{Perplexity (PPL)}
\begin{lstlisting}[label={lst:perplexity}]
import re
from math import exp

def extract_common_text(full_response: str) -> str:
    """
    Extracts only the shared text segments (Description & Decision) from the model output,
    filtering out Counterfactual (CF) Q&A sections to ensure fair comparison.
    """
    # 1. Extract all tagged sections (e.g., <accident>...</accident>)
    sections = re.findall(r'<[^>]+>.*?</[^>]+>', full_response)
    
    # 2. Extract prefix numbers (Accident Occurrence, Ego Involvement)
    prefix_match = re.match(r'^(\d+,\s*\d+)', full_response)
    prefix_numbers = prefix_match.group(1) if prefix_match else ""

    # 3. Filter common sections
    if len(sections) >= 5: 
        # Case: Model with CF reasoning (Output includes CF Q&A)
        # Format: [Desc1, Desc2, CF_Q, CF_A, Decision]
        # Action: Remove CF_Q (idx 2) and CF_A (idx 3)
        common_sections = [sections[0], sections[1], sections[-1]]
    else:
        # Case: Baseline Model (No CF Q&A)
        # Format: [Desc1, Desc2, Decision]
        # Action: Keep all sections
        common_sections = sections

    # 4. Reconstruct the common text
    return prefix_numbers + ", " + ", ".join(common_sections)


def compute_perplexity(model, video, prompt, target_text):
    """
    Computes the perplexity of the target_text given the video and prompt.
    """
    # 1. Construct full input sequence
    full_input = prompt + target_text
    
    # 2. Forward pass (Video + Text)
    # Inputs are tokenized and fed into the model to obtain logits
    logits = model(video, full_input)
    
    # 3. Set loss calculation range (Shifted Labels)
    # Exclude the prompt portion from loss computation; 
    # focus only on 'target_text'
    start_idx = len(tokenized_prompt)
    target_logits = logits[:, start_idx-1 : -1]
    target_labels = full_input_ids[:, start_idx:]
    
    # 4. Compute Cross Entropy Loss and Perplexity
    loss = CrossEntropyLoss(target_logits, target_labels)
    perplexity = exp(loss)
    
    return perplexity


# Main Evaluation Loop
def main(model_path, dataset_path):
    # Load model and dataset
    model = load_model(model_path)
    dataset = load_json(dataset_path)
    
    # Dictionary to store perplexity results
    ppl_results = {
        'overall': [],
        'stop': [],
        'decelerate': [],
        'keep': []
    }

    print("Starting Fair Perplexity Evaluation (Common Text Only)...")

    for sample in dataset:
        video_path = sample['video']
        # Original Ground Truth
        full_gt_text = sample['ground_truth']
        
        # Extract common text for fair comparison
        # Ensures that both CF and No-CF models 
        # are evaluated on identical text segments.
        common_target_text = extract_common_text(full_gt_text)
        
        # Extract decision class for breakdown analysis 
        # (Stop/Decelerate/Keep)
        decision_class = extract_decision_tag(full_gt_text)

        # Compute PPL based on the common text
        ppl = compute_perplexity(model, video_path, PROMPT, common_target_text)
        
        # Store results
        ppl_results['overall'].append(ppl)
        if decision_class in ppl_results:
            ppl_results[decision_class].append(ppl)

    # Compute and save final statistics
    print_statistics(ppl_results)
    save_results(ppl_results)

if __name__ == "__main__":
    main()
\end{lstlisting}



% \clearpage
% \subsection{Token Frequency Analysis}
% \begin{lstlisting}[label={lst:token_frequency}]
% SAFETY_TOKENS = {
%     "Longitudinal": 
%         {"stop", "brake", "decelerate", "halt", "slow", "yield"},
%     "Lateral":      
%         {"swerve", "steer", "turn", "change", "merge", "adjacent"},
%     "Event":        
%         {"collision", "collide", "crash", "accident"},
%     "State":        
%         {"safe", "unsafe", "danger", "risk", "caution"}
% }

% def analyze_safety_impact(dataset_base, dataset_cf):
%     """
%     Input:
%       dataset_base: Outputs from model without CF reasoning
%       dataset_cf:   Outputs from model with CF reasoning
%     Output:
%       coverage_rate: % of samples containing safety tokens
%       amplification: Ratio of token frequency (CF / Base)
%     """
    
%     freq = {"base": Counter(), "cf": Counter()}
%     coverage = {"base": Counter(), "cf": Counter()}
    
%     for sample_b, sample_c in zip(dataset_base, dataset_cf):
        
%         # Extract text (Description + QA/Reasoning, Counterfactuals)
%         text_b = extract_generated_text(sample_b) 
%         text_c = extract_generated_text(sample_c) 
        
%         tokens_b = tokenize(text_b)
%         tokens_c = tokenize(text_c)
        
%         for category, keywords in SAFETY_TOKENS.items():
%             count_b = sum(1 for t in tokens_b if t in keywords)
%             count_c = sum(1 for t in tokens_c if t in keywords)
            
%             # Update Total Frequency
%             freq["base"][category] += count_b
%             freq["cf"][category]   += count_c
            
%             # Update Sample Coverage (Binary hit)
%             if count_b > 0: coverage["base"][category] += 1
%             if count_c > 0: coverage["cf"][category]   += 1

%     results = {}
%     N = len(dataset_cf)
    
%     for cat in SAFETY_TOKENS:
%         rate_base = (coverage["base"][cat] / N) * 100
%         rate_cf   = (coverage["cf"][cat]   / N) * 100
        
%         f_b = freq["base"][cat]
%         f_c = freq["cf"][cat]
%         amp_factor = (f_c / f_b) if f_b > 0 else INF
        
%         results[cat] = {
%             "coverage_shift": f"{rate_base:.1f}% -> {rate_cf:.1f}%",
%             "amplification":  f"{amp_factor:.2f}x"
%         }
        
%     return results
% \end{lstlisting}


\clearpage
\section{Detailed Experimental Results}
\subsection{Supervised Fine-Tuning}
\label{appendix:passk_sft}

\begin{figure}[h]
    \centering
    \includegraphics[width=\columnwidth]{figure/passk/task_pass_at_k_grid_sft.pdf}
    \caption{Class-wise pass@$k$ performance across model scales and CF reasoning settings in SFT.}
    \label{fig:passk_sft}
\end{figure}

\paragraph{Analysis of Accident Recognition and Ego Inclusion.} 
As illustrated in Fig. \ref{fig:passk_sft}, CF reasoning consistently enhances the pass@1 performance for accident recognition and ego inclusion across all model scales. In smaller models, such as \texttt{InternVL2-1B} and \texttt{2B}, the performance gap between models with and without CF reasoning progressively diminishes as $k$ approaches 16, eventually converging to a similar baseline. However, a distinct trade-off emerges as the model size increases. In larger models like \texttt{Qwen2.5-VL-3B} and \texttt{VideoLLaMA2-7B}, while CF reasoning significantly boosts performance on accident-related labels ("accident", "ego involved"), it simultaneously leads to a performance degradation on non-accident labels ("no accident", "not involved") compared to the baseline. This suggests that CF reasoning does not merely provide additional information but fundamentally recalibrates the model's perception, making it inherently more sensitive to safety-critical risk factors.

\paragraph{Persistence of Safety-bias in Decision Fields.} 
This trend is further pronounced in the decision field, revealing a clear disparity in "confidence levels" tied to model scale. For the \texttt{Stop}, while smaller models show converging performance curves as k increases, larger models—specifically \texttt{Qwen2.5-VL-3B} and \texttt{VideoLLaMA2-7B}—maintain a persistent $\Delta$pass@$k$ that does not dissipate even at $k$=16. Conversely, for the \texttt{Dec.} and \texttt{Keep}, the baseline models without CF reasoning exhibit superior performance in these larger scales. This indicates that larger models internalize accident-related cues learned via CF reasoning as a robust "world prior." Consequently, in large-scale VLM, CF reasoning induces a deep-seated action distribution shift that remains resistant to stochastic sampling, acting as a dominant mechanism that drives the model toward conservative decision-making.


\clearpage
\subsection{Reinforcement Fine-Tuning}
\label{appendix:passk_rft}
\begin{figure}[t]
    \centering
    \includegraphics[width=\columnwidth]{figure/passk/task_pass_at_k_grid_rft.pdf}
    \caption{Class-wise pass@$k$ performance across model scales and CF reasoning settings in RFT.}
    \label{fig:passk_rft}
\end{figure}




% Δ Pass@k : w/ CF - w/o CF
\begin{figure}[h]
    \centering
    \includegraphics[width=\columnwidth]{figure/passk/delta_cf.pdf}
    \caption{$\Delta$pass@$k$ ($P_{\text{w/ CF}} - P_{\text{w/o CF}}$) for SFT and RFT across varying model scales and tasks, showing persistent gaps in larger models.}
    \label{fig:passk_delta}
\end{figure}


\paragraph{Distribution Shift and Performance Transfer in RFT.}
SFT models served as references for the RFT stage. As illustrated in Fig. \ref{fig:passk_rft}, accident-relevant tasks ("accident", "ego involved") exhibit a significant pass@1 distribution shift, mirroring SFT trends. While this gap converges at $k=16$ for \texttt{InternVL2.5-1B}, larger model like \texttt{Qwen2.5-VL-3B} demonstrate a persistent trade-off between accident sensitivity and non-accident performance.

\paragraph{Amplification of Decision-Making Bias via RFT.}
As shown in Fig. \ref{fig:passk_delta}, RFT amplifies the behavioral patterns established during SFT. For the \texttt{Stop} in \texttt{Qwen2.5-VL-3B}, the $\Delta$pass@$k$ gap between models with and without CF reasoning is notably wider than that observed in the SFT stage. Similarly, for \texttt{Dec.} and \texttt{Keep}, the initial gap at $k=1$ is more pronounced than in SFT, although differences diminish as $k$ approaches 16. This suggests RFT reinforces risk-related priors, leading to a more ingrained safety-bias.



\clearpage
\subsection{Attention Map}
\begin{figure}[h]
    \centering
    \includegraphics[width=\columnwidth]{figure/attention/00001_concat.jpg}
    \includegraphics[width=\columnwidth]{figure/attention/00009_concat.jpg}
    \includegraphics[width=\columnwidth]{figure/attention/00033_concat.jpg}
    \includegraphics[width=\columnwidth]{figure/attention/00038_concat.jpg}
    \includegraphics[width=\columnwidth]{figure/attention/00050_concat.jpg}
    \includegraphics[width=\columnwidth]{figure/attention/00067_concat.jpg}
    \includegraphics[width=\columnwidth]{figure/attention/00069_concat.jpg}
    \includegraphics[width=\columnwidth]{figure/attention/00101_concat.jpg}
    \includegraphics[width=\columnwidth]{figure/attention/00159_concat.jpg}
    \caption{Qualitative comparison of interference activations among SFT w/o CF, SFT w/ CF, RFT w/o CF, and RFT w/ CF based on Qwen2.5-VL-3B. Visualization of activation maps for causal factors immediately before accidents shows no noticeable differences across the models.}
    \label{fig:attention_map}
\end{figure}

%%%%%%%%%%%%%%%%%%%%%%%%%%%%%%%%%%%%%%%%%%%%%%%%%%%%%%%%%%%%

% \newpage
% \input{checklist.tex}


\end{document}